\documentclass[aspectratio=169,10pt]{beamer}
\makeatletter\def\input@path{{../../estilo/}{../estilo/}{estilo/}}\makeatother
\usepackage{devmobile}
\modulo[Flutter]{13}{Recursos do Dispositivo e IoT}{Permissões, câmera, QR code, GPS e mapas, sensores, biometria, NFC, notificações, segundo plano, arquivos, áudio, Bluetooth LE com ESP32, serial USB com Arduino e MQTT}

\begin{document}
\capa

\begin{frame}{Objetivos do módulo}
  \begin{conceito}[Ao final deste módulo você será capaz de]
    \begin{itemize}
      \item Pedir permissões corretamente no Android e no iOS.
      \item Usar câmera, galeria, leitor de QR code, GPS, mapas e sensores.
      \item Proteger ações com biometria e ler etiquetas NFC.
      \item Exibir notificações locais e executar tarefas em segundo plano.
      \item Compartilhar arquivos, abrir links, tocar e gravar áudio.
      \item Comunicar-se com placas ESP32 e Arduino por Bluetooth LE, USB serial e MQTT.
    \end{itemize}
  \end{conceito}
\end{frame}

\sumario

% =============================================================================
%  Fundamentos compartilhados --- Módulo 13: hardware, sensores e comunicação
%  Incluído pelas trilhas Kotlin e Flutter com % =============================================================================
%  Fundamentos compartilhados --- Módulo 13: hardware, sensores e comunicação
%  Incluído pelas trilhas Kotlin e Flutter com % =============================================================================
%  Fundamentos compartilhados --- Módulo 13: hardware, sensores e comunicação
%  Incluído pelas trilhas Kotlin e Flutter com \input{../fundamentos/f13-dispositivo}
% =============================================================================
\section{Fundamentos: hardware, sensores e comunicação}

\begin{frame}{Do mundo físico ao número}
  \centering
  \begin{tikzpicture}[node distance=4mm,every node/.style={font=\scriptsize}]
    \node[caixaambar,minimum width=2cm,minimum height=11mm] (g) {grandeza física\\ \tiny temperatura, aceleração};
    \node[caixa,minimum width=1.9cm,minimum height=11mm,right=of g] (t) {transdutor\\ \tiny vira tensão};
    \node[caixa,minimum width=2cm,minimum height=11mm,right=of t] (c) {condicionamento\\ \tiny amplifica, filtra};
    \node[caixapetroleo,minimum width=1.9cm,minimum height=11mm,right=of c] (a) {ADC\\ \tiny amostra e quantiza};
    \node[caixaescura,minimum width=2.2cm,minimum height=11mm,right=of a] (n) {número + unidade\\ \tiny no app};
    \draw[seta] (g) -- (t); \draw[seta] (t) -- (c); \draw[seta] (c) -- (a); \draw[seta] (a) -- (n);
  \end{tikzpicture}
  \raggedright
  \begin{columns}[T]
    \begin{column}{0.49\textwidth}
      \begin{definicao}[Conversor analógico-digital (ADC)]
        Mede a tensão em instantes regulares (\textbf{taxa de amostragem}) e a representa com um número de \textbf{bits}. Um ADC de 12 bits, como o do ESP32, distingue $2^{12} = 4096$ níveis.
      \end{definicao}
      \begin{pordentro}[Teorema da amostragem]
        Para reconstruir um sinal, amostre a mais que o dobro da maior frequência presente nele (Nyquist). Abaixo disso surgem frequências falsas (\emph{aliasing}).
      \end{pordentro}
    \end{column}
    \begin{column}{0.49\textwidth}
      \begin{definicao}[Sensores MEMS]
        Estruturas mecânicas microscópicas gravadas no silício. O \textbf{acelerômetro} mede aceleração (inclusive a da gravidade); o \textbf{giroscópio}, velocidade angular; o \textbf{magnetômetro}, o campo magnético; o \textbf{barômetro}, a pressão.
      \end{definicao}
      \begin{atencao}
        Toda medida tem ruído, deriva e erro de calibração: exiba com a precisão que o sensor realmente tem.
      \end{atencao}
    \end{column}
  \end{columns}
\end{frame}

\begin{frame}{Filtrando e combinando sensores}
  \begin{columns}[c]
    \begin{column}{0.48\textwidth}
      \centering
      \begin{tikzpicture}[x=0.55cm,y=0.9cm]
        \draw[->,cinzamedio] (0,0) -- (10.3,0) node[rotulo,right]{t};
        \draw[->,cinzamedio] (0,0) -- (0,2.2) node[rotulo,above]{°C};
        \draw[cinzamedio!70,line width=0.5pt] plot coordinates {(0,0.9) (0.5,1.3) (1,0.7) (1.5,1.2) (2,0.8) (2.5,1.4) (3,0.9) (3.5,1.1) (4,0.6) (4.5,1.3) (5,1.7) (5.5,1.4) (6,1.9) (6.5,1.5) (7,1.8) (7.5,1.4) (8,1.9) (8.5,1.6) (9,1.8) (9.5,1.5) (10,1.7)};
        \draw[ambar,line width=1.2pt] plot[smooth] coordinates {(0,0.95) (1,0.97) (2,0.98) (3,1.0) (4,0.98) (5,1.15) (6,1.4) (7,1.55) (8,1.64) (9,1.68) (10,1.68)};
        \node[rotulo,text=cinzamedio] at (2.5,1.75) {leitura bruta};
        \node[rotulo,text=ambar] at (7.6,1.0) {filtrada};
      \end{tikzpicture}
      \begin{definicao}[Média móvel exponencial]
        $y \leftarrow y + \alpha\,(x - y)$, com $\alpha$ entre 0 e 1. Pouco código, pouca memória: ótima para suavizar leituras.
      \end{definicao}
    \end{column}
    \begin{column}{0.5\textwidth}
      \begin{conceito}[Passa-baixas e passa-altas]
        Um filtro \textbf{passa-baixas} remove variações rápidas (ruído); um \textbf{passa-altas} remove variações lentas (deriva, gravidade). Quanto menor o $\alpha$, mais suave e mais atrasada fica a curva.
      \end{conceito}
      \begin{definicao}[Fusão de sensores]
        Combinar sensores para compensar os defeitos de cada um. O giroscópio é preciso no curto prazo e deriva; o acelerômetro é ruidoso, mas não deriva. Juntos, dão a orientação do aparelho.
      \end{definicao}
      \begin{dica}
        O sistema já entrega sensores fundidos, como o vetor de rotação. Prefira-os a reinventar a fusão.
      \end{dica}
    \end{column}
  \end{columns}
\end{frame}

\begin{frame}{Como o GPS descobre onde você está}
  \begin{columns}[c]
    \begin{column}{0.44\textwidth}
      \centering
      \begin{tikzpicture}
        \draw[azulpetroleo,line width=0.8pt] (0,0) circle (1.25);
        \draw[ambar,line width=0.8pt] (1.6,0.3) circle (1.05);
        \draw[cinzamedio,line width=0.8pt] (0.7,-1.3) circle (0.95);
        \fill[vermelhoerro] (0.95,-0.42) circle (2.2pt);
        \node[rotulo,text=vermelhoerro] at (1.45,-0.62) {você};
        \foreach \x/\y in {0/0,1.6/0.3,0.7/-1.3}{\fill[azulpetroleo] (\x,\y) circle (1.5pt);}
        \node[rotulo,text width=4cm] at (0.7,-2.6) {cada satélite define uma esfera de distância; a interseção é a posição};
      \end{tikzpicture}
    \end{column}
    \begin{column}{0.54\textwidth}
      \begin{definicao}[GNSS e trilateração]
        Satélites transmitem a hora de seus relógios atômicos; o receptor calcula a distância pelo tempo de viagem do sinal. Com \textbf{quatro} satélites, resolve latitude, longitude, altitude e o erro do próprio relógio.
      \end{definicao}
      \begin{pordentro}[Várias constelações e ajuda da rede]
        Celulares usam GPS (EUA), GLONASS (Rússia), Galileo (UE) e BeiDou (China). O A-GPS baixa pela internet a posição dos satélites e reduz o primeiro acerto de minutos para segundos.
      \end{pordentro}
      \begin{atencao}
        Dentro de prédios o sinal some e o sistema recorre a Wi-Fi, torres e sensores. Ao ar livre, a precisão típica é de alguns metros.
      \end{atencao}
    \end{column}
  \end{columns}
\end{frame}

\begin{frame}{Permissões e privacidade}
  \begin{columns}[T]
    \begin{column}{0.52\textwidth}
      \scriptsize
      \begin{tabularx}{\linewidth}{@{}l L L@{}}
        \toprule
        \cabecalho{Tipo} & \cabecalho{Concessão} & \cabecalho{Exemplos} \\
        \midrule
        Normal & automática, declarada no manifesto & internet, vibração \\
        Em tempo de execução & o usuário decide num diálogo & câmera, localização, Bluetooth, notificações \\
        Especial & tela própria nas configurações & alarmes exatos, sobrepor apps \\
        \bottomrule
      \end{tabularx}
      \normalsize
      \begin{historia}
        O Android pede permissões em tempo de execução desde a versão 6.0 (2015). Vieram depois a permissão de uma só vez (11), a localização aproximada (12) e o seletor de fotos sem permissão (13).
      \end{historia}
    \end{column}
    \begin{column}{0.46\textwidth}
      \begin{definicao}[Menor privilégio e minimização]
        Peça só o que o recurso exige, só quando ele for usado, e colete só os dados necessários. É boa engenharia e é exigência da LGPD (Módulo 15).
      \end{definicao}
      \begin{conceito}[No iOS]
        Cada permissão precisa de um texto de justificativa no \cod{Info.plist}; sem ele, o app é encerrado ao pedir o recurso.
      \end{conceito}
      \begin{dica}
        Um ponto colorido na barra de status avisa quando câmera ou microfone estão em uso, no Android e no iOS.
      \end{dica}
    \end{column}
  \end{columns}
\end{frame}

\begin{frame}{Bluetooth Low Energy por dentro}
  \begin{columns}[c]
    \begin{column}{0.46\textwidth}
      \centering
      \begin{tikzpicture}[every node/.style={font=\tiny}]
        \foreach \y/\n/\d/\estilo in {3/GAP/{papéis, anúncio e conexão}/caixaambar,2.35/GATT/{serviços e características}/caixaescura,1.7/ATT/{atributos com \emph{handle} e UUID}/caixapetroleo,1.05/L2CAP/{canais e fragmentação}/caixa,0.4/Enlace/{anúncios e eventos de conexão}/caixa,-0.25/Física/{2,4\,GHz, 40 canais}/caixa}{
          \node[\estilo,minimum width=1.3cm,minimum height=5.5mm,font=\tiny\bfseries] at (0,\y) {\n};
          \node[rotulo,anchor=west] at (0.75,\y) {\d};
        }
      \end{tikzpicture}
      \begin{historia}
        O BLE surgiu no Bluetooth 4.0 (2010) para sensores alimentados por baterias pequenas. O Bluetooth 5 (2016) trouxe o dobro de velocidade e um modo de longo alcance.
      \end{historia}
    \end{column}
    \begin{column}{0.52\textwidth}
      \begin{pordentro}[Por que gasta tão pouco]
        O rádio fica quase sempre desligado. O periférico anuncia em 3 dos 40 canais e, conectado, acorda só a cada \textbf{intervalo de conexão} (de 7,5\,ms a 4\,s) para trocar alguns bytes.
      \end{pordentro}
      \begin{definicao}[MTU]
        Tamanho máximo de um pacote ATT: 23 bytes por padrão (20 de dados úteis). Central e periférico podem negociar até 517 bytes.
      \end{definicao}
      \begin{dica}
        Leituras frequentes: use \textbf{notify}. Mais dados por segundo: aumente o MTU e reduza o intervalo de conexão, sabendo que isso custa bateria.
      \end{dica}
    \end{column}
  \end{columns}
\end{frame}

\begin{frame}{MQTT por dentro}
  \begin{columns}[c]
    \begin{column}{0.46\textwidth}
      \centering
      \begin{tikzpicture}[x=1cm,y=0.5cm,every node/.style={font=\tiny}]
        \node[rotulo] at (0,0.7) {\textbf{ESP32}};
        \node[rotulo] at (3.2,0.7) {\textbf{broker}};
        \draw[cinzamedio] (0,0.3) -- (0,-5.6); \draw[cinzamedio] (3.2,0.3) -- (3.2,-5.6);
        \node[rotulo,anchor=west] at (3.3,-0.1) {QoS 1};
        \draw[seta] (0,-0.2) -- node[rotulo,above,sloped]{PUBLISH} (3.2,-0.8);
        \draw[setaambar] (3.2,-1.1) -- node[rotulo,above,sloped]{PUBACK} (0,-1.7);
        \node[rotulo,anchor=west] at (3.3,-2.3) {QoS 2};
        \draw[seta] (0,-2.4) -- node[rotulo,above,sloped]{PUBLISH} (3.2,-3.0);
        \draw[setaambar] (3.2,-3.2) -- node[rotulo,above,sloped]{PUBREC} (0,-3.8);
        \draw[seta] (0,-4.0) -- node[rotulo,above,sloped]{PUBREL} (3.2,-4.6);
        \draw[setaambar] (3.2,-4.8) -- node[rotulo,above,sloped]{PUBCOMP} (0,-5.4);
      \end{tikzpicture}
    \end{column}
    \begin{column}{0.52\textwidth}
      \begin{historia}
        Andy Stanford-Clark (IBM) e Arlen Nipper criaram o MQTT em 1999 para monitorar oleodutos via satélite. Virou padrão OASIS em 2014; a versão 5.0 é de 2019.
      \end{historia}
      \begin{pordentro}[Leve de propósito]
        Roda sobre TCP (porta 1883, ou 8883 com TLS) e o cabeçalho fixo tem só 2 bytes. Mensagens de \emph{keep-alive} (PINGREQ e PINGRESP) mostram ao broker que a placa continua viva.
      \end{pordentro}
      \begin{conceito}[Sessão persistente]
        Com sessão persistente, o broker guarda as assinaturas e as mensagens de QoS 1 e 2 enquanto o cliente está desconectado e as entrega na volta.
      \end{conceito}
    \end{column}
  \end{columns}
\end{frame}

\begin{frame}{Comunicação serial e UART}
  \begin{columns}[c]
    \begin{column}{0.5\textwidth}
      \centering
      \begin{tikzpicture}[x=0.5cm,y=0.6cm]
        \draw[cinzamedio,dashed] (0,1) -- (14,1);
        \draw[azulpetroleo,line width=1pt] (0,1) -- (1,1) -- (1,0) -- (2,0) -- (2,1) -- (3,1) -- (3,0) -- (8,0) -- (8,1) -- (9,1) -- (9,0) -- (10,0) -- (10,1) -- (14,1);
        \node[rotulo] at (1.5,-0.45) {início};
        \foreach \b [count=\i from 0] in {1,0,0,0,0,0,1,0}{\node[rotulo] at (2.5+\i,1.35) {\b};}
        \node[rotulo] at (6,-0.45) {letra A (0x41), bit menos significativo primeiro};
        \node[rotulo] at (10.5,-0.45) {fim};
        \node[rotulo] at (12.6,0.6) {repouso};
      \end{tikzpicture}
      \begin{definicao}[UART]
        Comunicação serial \textbf{assíncrona}: não há fio de relógio. Os dois lados combinam a velocidade (\emph{baud rate}, como 115\,200) e o formato, quase sempre \textbf{8N1}: 8 bits de dados, sem paridade, 1 bit de parada.
      \end{definicao}
    \end{column}
    \begin{column}{0.48\textwidth}
      \begin{conceito}[Ligação]
        TX de um lado no RX do outro, e terra comum. Atenção aos níveis: o ESP32 trabalha em 3,3\,V, o Arduino Uno em 5\,V e o RS-232 em $\pm$12\,V.
      \end{conceito}
      \begin{pordentro}[Do USB à serial]
        A placa aparece no celular como dispositivo USB; um conversor (CH340, CP2102, FTDI) ou o próprio microcontrolador traduz para UART.
      \end{pordentro}
      \begin{dica}[Monte um protocolo]
        A serial só entrega bytes. Defina mensagens com delimitador (como \cod{\textbackslash n}) e, se o ambiente for ruidoso, uma soma de verificação.
      \end{dica}
    \end{column}
  \end{columns}
\end{frame}

\begin{frame}{NFC: comunicação por proximidade}
  \begin{columns}[T]
    \begin{column}{0.5\textwidth}
      \begin{definicao}[NFC]
        Comunicação a 13,56\,MHz por acoplamento indutivo, com alcance de poucos centímetros. A etiqueta passiva não tem bateria: é alimentada pelo campo do próprio leitor.
      \end{definicao}
      \scriptsize
      \begin{tabularx}{\linewidth}{@{}l L@{}}
        \toprule
        \cabecalho{Modo} & \cabecalho{Uso} \\
        \midrule
        Leitura e escrita & etiquetas com URL, texto ou identificador \\
        Emulação de cartão & o celular se passa por cartão (pagamentos, crachás) \\
        \bottomrule
      \end{tabularx}
    \end{column}
    \begin{column}{0.48\textwidth}
      \begin{definicao}[NDEF]
        Formato padrão das mensagens gravadas em etiquetas: uma lista de registros, cada um com tipo (URI, texto, MIME) e conteúdo.
      \end{definicao}
      \begin{exemplo}[No LabMonitor]
        Uma etiqueta colada na bancada guarda \cod{https://labmonitor.app/bancada/b3}: encostar o celular abre a tela da bancada por link profundo (Módulo 07).
      \end{exemplo}
      \begin{atencao}
        A distância curta é uma proteção, mas não substitui a autenticação: qualquer um pode ler ou copiar uma etiqueta sem proteção.
      \end{atencao}
    \end{column}
  \end{columns}
\end{frame}

\begin{frame}{Energia: o recurso mais escasso}
  \scriptsize
  \begin{tabularx}{\textwidth}{@{}l L L@{}}
    \toprule
    \cabecalho{Consumidor} & \cabecalho{Por que gasta} & \cabecalho{Como economizar} \\
    \midrule
    Tela & maior consumo do aparelho & tema escuro em telas OLED; nada de manter a tela ligada sem motivo \\
    GPS & receptor ligado e cálculos contínuos & localização aproximada, intervalos maiores, cercas virtuais \\
    Rede móvel & rádio de alta potência e tempo de cauda & agrupar requisições, \emph{push} em vez de \emph{polling} \\
    CPU & despertares frequentes & agendadores do sistema em vez de temporizadores próprios \\
    Bluetooth & varreduras longas & filtros por serviço e varreduras curtas \\
    \bottomrule
  \end{tabularx}
  \normalsize
  \vspace{4pt}
  \begin{columns}[T]
    \begin{column}{0.49\textwidth}
      \begin{historia}
        O Android criou o modo Doze (6.0, 2015), limitou serviços em segundo plano (8.0, 2017) e passou a classificar apps por frequência de uso (9.0, 2018). O iOS sempre restringiu o segundo plano a modos específicos.
      \end{historia}
    \end{column}
    \begin{column}{0.49\textwidth}
      \begin{conceito}[Regra geral]
        Deixe o sistema decidir \textbf{quando} executar: WorkManager no Android, \cod{BGTaskScheduler} no iOS. Ele agrupa o trabalho de vários apps e acorda o aparelho menos vezes.
      \end{conceito}
    \end{column}
  \end{columns}
\end{frame}

% =============================================================================
\section{Fundamentos: hardware, sensores e comunicação}

\begin{frame}{Do mundo físico ao número}
  \centering
  \begin{tikzpicture}[node distance=4mm,every node/.style={font=\scriptsize}]
    \node[caixaambar,minimum width=2cm,minimum height=11mm] (g) {grandeza física\\ \tiny temperatura, aceleração};
    \node[caixa,minimum width=1.9cm,minimum height=11mm,right=of g] (t) {transdutor\\ \tiny vira tensão};
    \node[caixa,minimum width=2cm,minimum height=11mm,right=of t] (c) {condicionamento\\ \tiny amplifica, filtra};
    \node[caixapetroleo,minimum width=1.9cm,minimum height=11mm,right=of c] (a) {ADC\\ \tiny amostra e quantiza};
    \node[caixaescura,minimum width=2.2cm,minimum height=11mm,right=of a] (n) {número + unidade\\ \tiny no app};
    \draw[seta] (g) -- (t); \draw[seta] (t) -- (c); \draw[seta] (c) -- (a); \draw[seta] (a) -- (n);
  \end{tikzpicture}
  \raggedright
  \begin{columns}[T]
    \begin{column}{0.49\textwidth}
      \begin{definicao}[Conversor analógico-digital (ADC)]
        Mede a tensão em instantes regulares (\textbf{taxa de amostragem}) e a representa com um número de \textbf{bits}. Um ADC de 12 bits, como o do ESP32, distingue $2^{12} = 4096$ níveis.
      \end{definicao}
      \begin{pordentro}[Teorema da amostragem]
        Para reconstruir um sinal, amostre a mais que o dobro da maior frequência presente nele (Nyquist). Abaixo disso surgem frequências falsas (\emph{aliasing}).
      \end{pordentro}
    \end{column}
    \begin{column}{0.49\textwidth}
      \begin{definicao}[Sensores MEMS]
        Estruturas mecânicas microscópicas gravadas no silício. O \textbf{acelerômetro} mede aceleração (inclusive a da gravidade); o \textbf{giroscópio}, velocidade angular; o \textbf{magnetômetro}, o campo magnético; o \textbf{barômetro}, a pressão.
      \end{definicao}
      \begin{atencao}
        Toda medida tem ruído, deriva e erro de calibração: exiba com a precisão que o sensor realmente tem.
      \end{atencao}
    \end{column}
  \end{columns}
\end{frame}

\begin{frame}{Filtrando e combinando sensores}
  \begin{columns}[c]
    \begin{column}{0.48\textwidth}
      \centering
      \begin{tikzpicture}[x=0.55cm,y=0.9cm]
        \draw[->,cinzamedio] (0,0) -- (10.3,0) node[rotulo,right]{t};
        \draw[->,cinzamedio] (0,0) -- (0,2.2) node[rotulo,above]{°C};
        \draw[cinzamedio!70,line width=0.5pt] plot coordinates {(0,0.9) (0.5,1.3) (1,0.7) (1.5,1.2) (2,0.8) (2.5,1.4) (3,0.9) (3.5,1.1) (4,0.6) (4.5,1.3) (5,1.7) (5.5,1.4) (6,1.9) (6.5,1.5) (7,1.8) (7.5,1.4) (8,1.9) (8.5,1.6) (9,1.8) (9.5,1.5) (10,1.7)};
        \draw[ambar,line width=1.2pt] plot[smooth] coordinates {(0,0.95) (1,0.97) (2,0.98) (3,1.0) (4,0.98) (5,1.15) (6,1.4) (7,1.55) (8,1.64) (9,1.68) (10,1.68)};
        \node[rotulo,text=cinzamedio] at (2.5,1.75) {leitura bruta};
        \node[rotulo,text=ambar] at (7.6,1.0) {filtrada};
      \end{tikzpicture}
      \begin{definicao}[Média móvel exponencial]
        $y \leftarrow y + \alpha\,(x - y)$, com $\alpha$ entre 0 e 1. Pouco código, pouca memória: ótima para suavizar leituras.
      \end{definicao}
    \end{column}
    \begin{column}{0.5\textwidth}
      \begin{conceito}[Passa-baixas e passa-altas]
        Um filtro \textbf{passa-baixas} remove variações rápidas (ruído); um \textbf{passa-altas} remove variações lentas (deriva, gravidade). Quanto menor o $\alpha$, mais suave e mais atrasada fica a curva.
      \end{conceito}
      \begin{definicao}[Fusão de sensores]
        Combinar sensores para compensar os defeitos de cada um. O giroscópio é preciso no curto prazo e deriva; o acelerômetro é ruidoso, mas não deriva. Juntos, dão a orientação do aparelho.
      \end{definicao}
      \begin{dica}
        O sistema já entrega sensores fundidos, como o vetor de rotação. Prefira-os a reinventar a fusão.
      \end{dica}
    \end{column}
  \end{columns}
\end{frame}

\begin{frame}{Como o GPS descobre onde você está}
  \begin{columns}[c]
    \begin{column}{0.44\textwidth}
      \centering
      \begin{tikzpicture}
        \draw[azulpetroleo,line width=0.8pt] (0,0) circle (1.25);
        \draw[ambar,line width=0.8pt] (1.6,0.3) circle (1.05);
        \draw[cinzamedio,line width=0.8pt] (0.7,-1.3) circle (0.95);
        \fill[vermelhoerro] (0.95,-0.42) circle (2.2pt);
        \node[rotulo,text=vermelhoerro] at (1.45,-0.62) {você};
        \foreach \x/\y in {0/0,1.6/0.3,0.7/-1.3}{\fill[azulpetroleo] (\x,\y) circle (1.5pt);}
        \node[rotulo,text width=4cm] at (0.7,-2.6) {cada satélite define uma esfera de distância; a interseção é a posição};
      \end{tikzpicture}
    \end{column}
    \begin{column}{0.54\textwidth}
      \begin{definicao}[GNSS e trilateração]
        Satélites transmitem a hora de seus relógios atômicos; o receptor calcula a distância pelo tempo de viagem do sinal. Com \textbf{quatro} satélites, resolve latitude, longitude, altitude e o erro do próprio relógio.
      \end{definicao}
      \begin{pordentro}[Várias constelações e ajuda da rede]
        Celulares usam GPS (EUA), GLONASS (Rússia), Galileo (UE) e BeiDou (China). O A-GPS baixa pela internet a posição dos satélites e reduz o primeiro acerto de minutos para segundos.
      \end{pordentro}
      \begin{atencao}
        Dentro de prédios o sinal some e o sistema recorre a Wi-Fi, torres e sensores. Ao ar livre, a precisão típica é de alguns metros.
      \end{atencao}
    \end{column}
  \end{columns}
\end{frame}

\begin{frame}{Permissões e privacidade}
  \begin{columns}[T]
    \begin{column}{0.52\textwidth}
      \scriptsize
      \begin{tabularx}{\linewidth}{@{}l L L@{}}
        \toprule
        \cabecalho{Tipo} & \cabecalho{Concessão} & \cabecalho{Exemplos} \\
        \midrule
        Normal & automática, declarada no manifesto & internet, vibração \\
        Em tempo de execução & o usuário decide num diálogo & câmera, localização, Bluetooth, notificações \\
        Especial & tela própria nas configurações & alarmes exatos, sobrepor apps \\
        \bottomrule
      \end{tabularx}
      \normalsize
      \begin{historia}
        O Android pede permissões em tempo de execução desde a versão 6.0 (2015). Vieram depois a permissão de uma só vez (11), a localização aproximada (12) e o seletor de fotos sem permissão (13).
      \end{historia}
    \end{column}
    \begin{column}{0.46\textwidth}
      \begin{definicao}[Menor privilégio e minimização]
        Peça só o que o recurso exige, só quando ele for usado, e colete só os dados necessários. É boa engenharia e é exigência da LGPD (Módulo 15).
      \end{definicao}
      \begin{conceito}[No iOS]
        Cada permissão precisa de um texto de justificativa no \cod{Info.plist}; sem ele, o app é encerrado ao pedir o recurso.
      \end{conceito}
      \begin{dica}
        Um ponto colorido na barra de status avisa quando câmera ou microfone estão em uso, no Android e no iOS.
      \end{dica}
    \end{column}
  \end{columns}
\end{frame}

\begin{frame}{Bluetooth Low Energy por dentro}
  \begin{columns}[c]
    \begin{column}{0.46\textwidth}
      \centering
      \begin{tikzpicture}[every node/.style={font=\tiny}]
        \foreach \y/\n/\d/\estilo in {3/GAP/{papéis, anúncio e conexão}/caixaambar,2.35/GATT/{serviços e características}/caixaescura,1.7/ATT/{atributos com \emph{handle} e UUID}/caixapetroleo,1.05/L2CAP/{canais e fragmentação}/caixa,0.4/Enlace/{anúncios e eventos de conexão}/caixa,-0.25/Física/{2,4\,GHz, 40 canais}/caixa}{
          \node[\estilo,minimum width=1.3cm,minimum height=5.5mm,font=\tiny\bfseries] at (0,\y) {\n};
          \node[rotulo,anchor=west] at (0.75,\y) {\d};
        }
      \end{tikzpicture}
      \begin{historia}
        O BLE surgiu no Bluetooth 4.0 (2010) para sensores alimentados por baterias pequenas. O Bluetooth 5 (2016) trouxe o dobro de velocidade e um modo de longo alcance.
      \end{historia}
    \end{column}
    \begin{column}{0.52\textwidth}
      \begin{pordentro}[Por que gasta tão pouco]
        O rádio fica quase sempre desligado. O periférico anuncia em 3 dos 40 canais e, conectado, acorda só a cada \textbf{intervalo de conexão} (de 7,5\,ms a 4\,s) para trocar alguns bytes.
      \end{pordentro}
      \begin{definicao}[MTU]
        Tamanho máximo de um pacote ATT: 23 bytes por padrão (20 de dados úteis). Central e periférico podem negociar até 517 bytes.
      \end{definicao}
      \begin{dica}
        Leituras frequentes: use \textbf{notify}. Mais dados por segundo: aumente o MTU e reduza o intervalo de conexão, sabendo que isso custa bateria.
      \end{dica}
    \end{column}
  \end{columns}
\end{frame}

\begin{frame}{MQTT por dentro}
  \begin{columns}[c]
    \begin{column}{0.46\textwidth}
      \centering
      \begin{tikzpicture}[x=1cm,y=0.5cm,every node/.style={font=\tiny}]
        \node[rotulo] at (0,0.7) {\textbf{ESP32}};
        \node[rotulo] at (3.2,0.7) {\textbf{broker}};
        \draw[cinzamedio] (0,0.3) -- (0,-5.6); \draw[cinzamedio] (3.2,0.3) -- (3.2,-5.6);
        \node[rotulo,anchor=west] at (3.3,-0.1) {QoS 1};
        \draw[seta] (0,-0.2) -- node[rotulo,above,sloped]{PUBLISH} (3.2,-0.8);
        \draw[setaambar] (3.2,-1.1) -- node[rotulo,above,sloped]{PUBACK} (0,-1.7);
        \node[rotulo,anchor=west] at (3.3,-2.3) {QoS 2};
        \draw[seta] (0,-2.4) -- node[rotulo,above,sloped]{PUBLISH} (3.2,-3.0);
        \draw[setaambar] (3.2,-3.2) -- node[rotulo,above,sloped]{PUBREC} (0,-3.8);
        \draw[seta] (0,-4.0) -- node[rotulo,above,sloped]{PUBREL} (3.2,-4.6);
        \draw[setaambar] (3.2,-4.8) -- node[rotulo,above,sloped]{PUBCOMP} (0,-5.4);
      \end{tikzpicture}
    \end{column}
    \begin{column}{0.52\textwidth}
      \begin{historia}
        Andy Stanford-Clark (IBM) e Arlen Nipper criaram o MQTT em 1999 para monitorar oleodutos via satélite. Virou padrão OASIS em 2014; a versão 5.0 é de 2019.
      \end{historia}
      \begin{pordentro}[Leve de propósito]
        Roda sobre TCP (porta 1883, ou 8883 com TLS) e o cabeçalho fixo tem só 2 bytes. Mensagens de \emph{keep-alive} (PINGREQ e PINGRESP) mostram ao broker que a placa continua viva.
      \end{pordentro}
      \begin{conceito}[Sessão persistente]
        Com sessão persistente, o broker guarda as assinaturas e as mensagens de QoS 1 e 2 enquanto o cliente está desconectado e as entrega na volta.
      \end{conceito}
    \end{column}
  \end{columns}
\end{frame}

\begin{frame}{Comunicação serial e UART}
  \begin{columns}[c]
    \begin{column}{0.5\textwidth}
      \centering
      \begin{tikzpicture}[x=0.5cm,y=0.6cm]
        \draw[cinzamedio,dashed] (0,1) -- (14,1);
        \draw[azulpetroleo,line width=1pt] (0,1) -- (1,1) -- (1,0) -- (2,0) -- (2,1) -- (3,1) -- (3,0) -- (8,0) -- (8,1) -- (9,1) -- (9,0) -- (10,0) -- (10,1) -- (14,1);
        \node[rotulo] at (1.5,-0.45) {início};
        \foreach \b [count=\i from 0] in {1,0,0,0,0,0,1,0}{\node[rotulo] at (2.5+\i,1.35) {\b};}
        \node[rotulo] at (6,-0.45) {letra A (0x41), bit menos significativo primeiro};
        \node[rotulo] at (10.5,-0.45) {fim};
        \node[rotulo] at (12.6,0.6) {repouso};
      \end{tikzpicture}
      \begin{definicao}[UART]
        Comunicação serial \textbf{assíncrona}: não há fio de relógio. Os dois lados combinam a velocidade (\emph{baud rate}, como 115\,200) e o formato, quase sempre \textbf{8N1}: 8 bits de dados, sem paridade, 1 bit de parada.
      \end{definicao}
    \end{column}
    \begin{column}{0.48\textwidth}
      \begin{conceito}[Ligação]
        TX de um lado no RX do outro, e terra comum. Atenção aos níveis: o ESP32 trabalha em 3,3\,V, o Arduino Uno em 5\,V e o RS-232 em $\pm$12\,V.
      \end{conceito}
      \begin{pordentro}[Do USB à serial]
        A placa aparece no celular como dispositivo USB; um conversor (CH340, CP2102, FTDI) ou o próprio microcontrolador traduz para UART.
      \end{pordentro}
      \begin{dica}[Monte um protocolo]
        A serial só entrega bytes. Defina mensagens com delimitador (como \cod{\textbackslash n}) e, se o ambiente for ruidoso, uma soma de verificação.
      \end{dica}
    \end{column}
  \end{columns}
\end{frame}

\begin{frame}{NFC: comunicação por proximidade}
  \begin{columns}[T]
    \begin{column}{0.5\textwidth}
      \begin{definicao}[NFC]
        Comunicação a 13,56\,MHz por acoplamento indutivo, com alcance de poucos centímetros. A etiqueta passiva não tem bateria: é alimentada pelo campo do próprio leitor.
      \end{definicao}
      \scriptsize
      \begin{tabularx}{\linewidth}{@{}l L@{}}
        \toprule
        \cabecalho{Modo} & \cabecalho{Uso} \\
        \midrule
        Leitura e escrita & etiquetas com URL, texto ou identificador \\
        Emulação de cartão & o celular se passa por cartão (pagamentos, crachás) \\
        \bottomrule
      \end{tabularx}
    \end{column}
    \begin{column}{0.48\textwidth}
      \begin{definicao}[NDEF]
        Formato padrão das mensagens gravadas em etiquetas: uma lista de registros, cada um com tipo (URI, texto, MIME) e conteúdo.
      \end{definicao}
      \begin{exemplo}[No LabMonitor]
        Uma etiqueta colada na bancada guarda \cod{https://labmonitor.app/bancada/b3}: encostar o celular abre a tela da bancada por link profundo (Módulo 07).
      \end{exemplo}
      \begin{atencao}
        A distância curta é uma proteção, mas não substitui a autenticação: qualquer um pode ler ou copiar uma etiqueta sem proteção.
      \end{atencao}
    \end{column}
  \end{columns}
\end{frame}

\begin{frame}{Energia: o recurso mais escasso}
  \scriptsize
  \begin{tabularx}{\textwidth}{@{}l L L@{}}
    \toprule
    \cabecalho{Consumidor} & \cabecalho{Por que gasta} & \cabecalho{Como economizar} \\
    \midrule
    Tela & maior consumo do aparelho & tema escuro em telas OLED; nada de manter a tela ligada sem motivo \\
    GPS & receptor ligado e cálculos contínuos & localização aproximada, intervalos maiores, cercas virtuais \\
    Rede móvel & rádio de alta potência e tempo de cauda & agrupar requisições, \emph{push} em vez de \emph{polling} \\
    CPU & despertares frequentes & agendadores do sistema em vez de temporizadores próprios \\
    Bluetooth & varreduras longas & filtros por serviço e varreduras curtas \\
    \bottomrule
  \end{tabularx}
  \normalsize
  \vspace{4pt}
  \begin{columns}[T]
    \begin{column}{0.49\textwidth}
      \begin{historia}
        O Android criou o modo Doze (6.0, 2015), limitou serviços em segundo plano (8.0, 2017) e passou a classificar apps por frequência de uso (9.0, 2018). O iOS sempre restringiu o segundo plano a modos específicos.
      \end{historia}
    \end{column}
    \begin{column}{0.49\textwidth}
      \begin{conceito}[Regra geral]
        Deixe o sistema decidir \textbf{quando} executar: WorkManager no Android, \cod{BGTaskScheduler} no iOS. Ele agrupa o trabalho de vários apps e acorda o aparelho menos vezes.
      \end{conceito}
    \end{column}
  \end{columns}
\end{frame}

% =============================================================================
\section{Fundamentos: hardware, sensores e comunicação}

\begin{frame}{Do mundo físico ao número}
  \centering
  \begin{tikzpicture}[node distance=4mm,every node/.style={font=\scriptsize}]
    \node[caixaambar,minimum width=2cm,minimum height=11mm] (g) {grandeza física\\ \tiny temperatura, aceleração};
    \node[caixa,minimum width=1.9cm,minimum height=11mm,right=of g] (t) {transdutor\\ \tiny vira tensão};
    \node[caixa,minimum width=2cm,minimum height=11mm,right=of t] (c) {condicionamento\\ \tiny amplifica, filtra};
    \node[caixapetroleo,minimum width=1.9cm,minimum height=11mm,right=of c] (a) {ADC\\ \tiny amostra e quantiza};
    \node[caixaescura,minimum width=2.2cm,minimum height=11mm,right=of a] (n) {número + unidade\\ \tiny no app};
    \draw[seta] (g) -- (t); \draw[seta] (t) -- (c); \draw[seta] (c) -- (a); \draw[seta] (a) -- (n);
  \end{tikzpicture}
  \raggedright
  \begin{columns}[T]
    \begin{column}{0.49\textwidth}
      \begin{definicao}[Conversor analógico-digital (ADC)]
        Mede a tensão em instantes regulares (\textbf{taxa de amostragem}) e a representa com um número de \textbf{bits}. Um ADC de 12 bits, como o do ESP32, distingue $2^{12} = 4096$ níveis.
      \end{definicao}
      \begin{pordentro}[Teorema da amostragem]
        Para reconstruir um sinal, amostre a mais que o dobro da maior frequência presente nele (Nyquist). Abaixo disso surgem frequências falsas (\emph{aliasing}).
      \end{pordentro}
    \end{column}
    \begin{column}{0.49\textwidth}
      \begin{definicao}[Sensores MEMS]
        Estruturas mecânicas microscópicas gravadas no silício. O \textbf{acelerômetro} mede aceleração (inclusive a da gravidade); o \textbf{giroscópio}, velocidade angular; o \textbf{magnetômetro}, o campo magnético; o \textbf{barômetro}, a pressão.
      \end{definicao}
      \begin{atencao}
        Toda medida tem ruído, deriva e erro de calibração: exiba com a precisão que o sensor realmente tem.
      \end{atencao}
    \end{column}
  \end{columns}
\end{frame}

\begin{frame}{Filtrando e combinando sensores}
  \begin{columns}[c]
    \begin{column}{0.48\textwidth}
      \centering
      \begin{tikzpicture}[x=0.55cm,y=0.9cm]
        \draw[->,cinzamedio] (0,0) -- (10.3,0) node[rotulo,right]{t};
        \draw[->,cinzamedio] (0,0) -- (0,2.2) node[rotulo,above]{°C};
        \draw[cinzamedio!70,line width=0.5pt] plot coordinates {(0,0.9) (0.5,1.3) (1,0.7) (1.5,1.2) (2,0.8) (2.5,1.4) (3,0.9) (3.5,1.1) (4,0.6) (4.5,1.3) (5,1.7) (5.5,1.4) (6,1.9) (6.5,1.5) (7,1.8) (7.5,1.4) (8,1.9) (8.5,1.6) (9,1.8) (9.5,1.5) (10,1.7)};
        \draw[ambar,line width=1.2pt] plot[smooth] coordinates {(0,0.95) (1,0.97) (2,0.98) (3,1.0) (4,0.98) (5,1.15) (6,1.4) (7,1.55) (8,1.64) (9,1.68) (10,1.68)};
        \node[rotulo,text=cinzamedio] at (2.5,1.75) {leitura bruta};
        \node[rotulo,text=ambar] at (7.6,1.0) {filtrada};
      \end{tikzpicture}
      \begin{definicao}[Média móvel exponencial]
        $y \leftarrow y + \alpha\,(x - y)$, com $\alpha$ entre 0 e 1. Pouco código, pouca memória: ótima para suavizar leituras.
      \end{definicao}
    \end{column}
    \begin{column}{0.5\textwidth}
      \begin{conceito}[Passa-baixas e passa-altas]
        Um filtro \textbf{passa-baixas} remove variações rápidas (ruído); um \textbf{passa-altas} remove variações lentas (deriva, gravidade). Quanto menor o $\alpha$, mais suave e mais atrasada fica a curva.
      \end{conceito}
      \begin{definicao}[Fusão de sensores]
        Combinar sensores para compensar os defeitos de cada um. O giroscópio é preciso no curto prazo e deriva; o acelerômetro é ruidoso, mas não deriva. Juntos, dão a orientação do aparelho.
      \end{definicao}
      \begin{dica}
        O sistema já entrega sensores fundidos, como o vetor de rotação. Prefira-os a reinventar a fusão.
      \end{dica}
    \end{column}
  \end{columns}
\end{frame}

\begin{frame}{Como o GPS descobre onde você está}
  \begin{columns}[c]
    \begin{column}{0.44\textwidth}
      \centering
      \begin{tikzpicture}
        \draw[azulpetroleo,line width=0.8pt] (0,0) circle (1.25);
        \draw[ambar,line width=0.8pt] (1.6,0.3) circle (1.05);
        \draw[cinzamedio,line width=0.8pt] (0.7,-1.3) circle (0.95);
        \fill[vermelhoerro] (0.95,-0.42) circle (2.2pt);
        \node[rotulo,text=vermelhoerro] at (1.45,-0.62) {você};
        \foreach \x/\y in {0/0,1.6/0.3,0.7/-1.3}{\fill[azulpetroleo] (\x,\y) circle (1.5pt);}
        \node[rotulo,text width=4cm] at (0.7,-2.6) {cada satélite define uma esfera de distância; a interseção é a posição};
      \end{tikzpicture}
    \end{column}
    \begin{column}{0.54\textwidth}
      \begin{definicao}[GNSS e trilateração]
        Satélites transmitem a hora de seus relógios atômicos; o receptor calcula a distância pelo tempo de viagem do sinal. Com \textbf{quatro} satélites, resolve latitude, longitude, altitude e o erro do próprio relógio.
      \end{definicao}
      \begin{pordentro}[Várias constelações e ajuda da rede]
        Celulares usam GPS (EUA), GLONASS (Rússia), Galileo (UE) e BeiDou (China). O A-GPS baixa pela internet a posição dos satélites e reduz o primeiro acerto de minutos para segundos.
      \end{pordentro}
      \begin{atencao}
        Dentro de prédios o sinal some e o sistema recorre a Wi-Fi, torres e sensores. Ao ar livre, a precisão típica é de alguns metros.
      \end{atencao}
    \end{column}
  \end{columns}
\end{frame}

\begin{frame}{Permissões e privacidade}
  \begin{columns}[T]
    \begin{column}{0.52\textwidth}
      \scriptsize
      \begin{tabularx}{\linewidth}{@{}l L L@{}}
        \toprule
        \cabecalho{Tipo} & \cabecalho{Concessão} & \cabecalho{Exemplos} \\
        \midrule
        Normal & automática, declarada no manifesto & internet, vibração \\
        Em tempo de execução & o usuário decide num diálogo & câmera, localização, Bluetooth, notificações \\
        Especial & tela própria nas configurações & alarmes exatos, sobrepor apps \\
        \bottomrule
      \end{tabularx}
      \normalsize
      \begin{historia}
        O Android pede permissões em tempo de execução desde a versão 6.0 (2015). Vieram depois a permissão de uma só vez (11), a localização aproximada (12) e o seletor de fotos sem permissão (13).
      \end{historia}
    \end{column}
    \begin{column}{0.46\textwidth}
      \begin{definicao}[Menor privilégio e minimização]
        Peça só o que o recurso exige, só quando ele for usado, e colete só os dados necessários. É boa engenharia e é exigência da LGPD (Módulo 15).
      \end{definicao}
      \begin{conceito}[No iOS]
        Cada permissão precisa de um texto de justificativa no \cod{Info.plist}; sem ele, o app é encerrado ao pedir o recurso.
      \end{conceito}
      \begin{dica}
        Um ponto colorido na barra de status avisa quando câmera ou microfone estão em uso, no Android e no iOS.
      \end{dica}
    \end{column}
  \end{columns}
\end{frame}

\begin{frame}{Bluetooth Low Energy por dentro}
  \begin{columns}[c]
    \begin{column}{0.46\textwidth}
      \centering
      \begin{tikzpicture}[every node/.style={font=\tiny}]
        \foreach \y/\n/\d/\estilo in {3/GAP/{papéis, anúncio e conexão}/caixaambar,2.35/GATT/{serviços e características}/caixaescura,1.7/ATT/{atributos com \emph{handle} e UUID}/caixapetroleo,1.05/L2CAP/{canais e fragmentação}/caixa,0.4/Enlace/{anúncios e eventos de conexão}/caixa,-0.25/Física/{2,4\,GHz, 40 canais}/caixa}{
          \node[\estilo,minimum width=1.3cm,minimum height=5.5mm,font=\tiny\bfseries] at (0,\y) {\n};
          \node[rotulo,anchor=west] at (0.75,\y) {\d};
        }
      \end{tikzpicture}
      \begin{historia}
        O BLE surgiu no Bluetooth 4.0 (2010) para sensores alimentados por baterias pequenas. O Bluetooth 5 (2016) trouxe o dobro de velocidade e um modo de longo alcance.
      \end{historia}
    \end{column}
    \begin{column}{0.52\textwidth}
      \begin{pordentro}[Por que gasta tão pouco]
        O rádio fica quase sempre desligado. O periférico anuncia em 3 dos 40 canais e, conectado, acorda só a cada \textbf{intervalo de conexão} (de 7,5\,ms a 4\,s) para trocar alguns bytes.
      \end{pordentro}
      \begin{definicao}[MTU]
        Tamanho máximo de um pacote ATT: 23 bytes por padrão (20 de dados úteis). Central e periférico podem negociar até 517 bytes.
      \end{definicao}
      \begin{dica}
        Leituras frequentes: use \textbf{notify}. Mais dados por segundo: aumente o MTU e reduza o intervalo de conexão, sabendo que isso custa bateria.
      \end{dica}
    \end{column}
  \end{columns}
\end{frame}

\begin{frame}{MQTT por dentro}
  \begin{columns}[c]
    \begin{column}{0.46\textwidth}
      \centering
      \begin{tikzpicture}[x=1cm,y=0.5cm,every node/.style={font=\tiny}]
        \node[rotulo] at (0,0.7) {\textbf{ESP32}};
        \node[rotulo] at (3.2,0.7) {\textbf{broker}};
        \draw[cinzamedio] (0,0.3) -- (0,-5.6); \draw[cinzamedio] (3.2,0.3) -- (3.2,-5.6);
        \node[rotulo,anchor=west] at (3.3,-0.1) {QoS 1};
        \draw[seta] (0,-0.2) -- node[rotulo,above,sloped]{PUBLISH} (3.2,-0.8);
        \draw[setaambar] (3.2,-1.1) -- node[rotulo,above,sloped]{PUBACK} (0,-1.7);
        \node[rotulo,anchor=west] at (3.3,-2.3) {QoS 2};
        \draw[seta] (0,-2.4) -- node[rotulo,above,sloped]{PUBLISH} (3.2,-3.0);
        \draw[setaambar] (3.2,-3.2) -- node[rotulo,above,sloped]{PUBREC} (0,-3.8);
        \draw[seta] (0,-4.0) -- node[rotulo,above,sloped]{PUBREL} (3.2,-4.6);
        \draw[setaambar] (3.2,-4.8) -- node[rotulo,above,sloped]{PUBCOMP} (0,-5.4);
      \end{tikzpicture}
    \end{column}
    \begin{column}{0.52\textwidth}
      \begin{historia}
        Andy Stanford-Clark (IBM) e Arlen Nipper criaram o MQTT em 1999 para monitorar oleodutos via satélite. Virou padrão OASIS em 2014; a versão 5.0 é de 2019.
      \end{historia}
      \begin{pordentro}[Leve de propósito]
        Roda sobre TCP (porta 1883, ou 8883 com TLS) e o cabeçalho fixo tem só 2 bytes. Mensagens de \emph{keep-alive} (PINGREQ e PINGRESP) mostram ao broker que a placa continua viva.
      \end{pordentro}
      \begin{conceito}[Sessão persistente]
        Com sessão persistente, o broker guarda as assinaturas e as mensagens de QoS 1 e 2 enquanto o cliente está desconectado e as entrega na volta.
      \end{conceito}
    \end{column}
  \end{columns}
\end{frame}

\begin{frame}{Comunicação serial e UART}
  \begin{columns}[c]
    \begin{column}{0.5\textwidth}
      \centering
      \begin{tikzpicture}[x=0.5cm,y=0.6cm]
        \draw[cinzamedio,dashed] (0,1) -- (14,1);
        \draw[azulpetroleo,line width=1pt] (0,1) -- (1,1) -- (1,0) -- (2,0) -- (2,1) -- (3,1) -- (3,0) -- (8,0) -- (8,1) -- (9,1) -- (9,0) -- (10,0) -- (10,1) -- (14,1);
        \node[rotulo] at (1.5,-0.45) {início};
        \foreach \b [count=\i from 0] in {1,0,0,0,0,0,1,0}{\node[rotulo] at (2.5+\i,1.35) {\b};}
        \node[rotulo] at (6,-0.45) {letra A (0x41), bit menos significativo primeiro};
        \node[rotulo] at (10.5,-0.45) {fim};
        \node[rotulo] at (12.6,0.6) {repouso};
      \end{tikzpicture}
      \begin{definicao}[UART]
        Comunicação serial \textbf{assíncrona}: não há fio de relógio. Os dois lados combinam a velocidade (\emph{baud rate}, como 115\,200) e o formato, quase sempre \textbf{8N1}: 8 bits de dados, sem paridade, 1 bit de parada.
      \end{definicao}
    \end{column}
    \begin{column}{0.48\textwidth}
      \begin{conceito}[Ligação]
        TX de um lado no RX do outro, e terra comum. Atenção aos níveis: o ESP32 trabalha em 3,3\,V, o Arduino Uno em 5\,V e o RS-232 em $\pm$12\,V.
      \end{conceito}
      \begin{pordentro}[Do USB à serial]
        A placa aparece no celular como dispositivo USB; um conversor (CH340, CP2102, FTDI) ou o próprio microcontrolador traduz para UART.
      \end{pordentro}
      \begin{dica}[Monte um protocolo]
        A serial só entrega bytes. Defina mensagens com delimitador (como \cod{\textbackslash n}) e, se o ambiente for ruidoso, uma soma de verificação.
      \end{dica}
    \end{column}
  \end{columns}
\end{frame}

\begin{frame}{NFC: comunicação por proximidade}
  \begin{columns}[T]
    \begin{column}{0.5\textwidth}
      \begin{definicao}[NFC]
        Comunicação a 13,56\,MHz por acoplamento indutivo, com alcance de poucos centímetros. A etiqueta passiva não tem bateria: é alimentada pelo campo do próprio leitor.
      \end{definicao}
      \scriptsize
      \begin{tabularx}{\linewidth}{@{}l L@{}}
        \toprule
        \cabecalho{Modo} & \cabecalho{Uso} \\
        \midrule
        Leitura e escrita & etiquetas com URL, texto ou identificador \\
        Emulação de cartão & o celular se passa por cartão (pagamentos, crachás) \\
        \bottomrule
      \end{tabularx}
    \end{column}
    \begin{column}{0.48\textwidth}
      \begin{definicao}[NDEF]
        Formato padrão das mensagens gravadas em etiquetas: uma lista de registros, cada um com tipo (URI, texto, MIME) e conteúdo.
      \end{definicao}
      \begin{exemplo}[No LabMonitor]
        Uma etiqueta colada na bancada guarda \cod{https://labmonitor.app/bancada/b3}: encostar o celular abre a tela da bancada por link profundo (Módulo 07).
      \end{exemplo}
      \begin{atencao}
        A distância curta é uma proteção, mas não substitui a autenticação: qualquer um pode ler ou copiar uma etiqueta sem proteção.
      \end{atencao}
    \end{column}
  \end{columns}
\end{frame}

\begin{frame}{Energia: o recurso mais escasso}
  \scriptsize
  \begin{tabularx}{\textwidth}{@{}l L L@{}}
    \toprule
    \cabecalho{Consumidor} & \cabecalho{Por que gasta} & \cabecalho{Como economizar} \\
    \midrule
    Tela & maior consumo do aparelho & tema escuro em telas OLED; nada de manter a tela ligada sem motivo \\
    GPS & receptor ligado e cálculos contínuos & localização aproximada, intervalos maiores, cercas virtuais \\
    Rede móvel & rádio de alta potência e tempo de cauda & agrupar requisições, \emph{push} em vez de \emph{polling} \\
    CPU & despertares frequentes & agendadores do sistema em vez de temporizadores próprios \\
    Bluetooth & varreduras longas & filtros por serviço e varreduras curtas \\
    \bottomrule
  \end{tabularx}
  \normalsize
  \vspace{4pt}
  \begin{columns}[T]
    \begin{column}{0.49\textwidth}
      \begin{historia}
        O Android criou o modo Doze (6.0, 2015), limitou serviços em segundo plano (8.0, 2017) e passou a classificar apps por frequência de uso (9.0, 2018). O iOS sempre restringiu o segundo plano a modos específicos.
      \end{historia}
    \end{column}
    \begin{column}{0.49\textwidth}
      \begin{conceito}[Regra geral]
        Deixe o sistema decidir \textbf{quando} executar: WorkManager no Android, \cod{BGTaskScheduler} no iOS. Ele agrupa o trabalho de vários apps e acorda o aparelho menos vezes.
      \end{conceito}
    \end{column}
  \end{columns}
\end{frame}


% =============================================================================
\section{Como o Flutter acessa o hardware}

\begin{frame}{Plugins: Dart por cima, nativo por baixo}
  \begin{columns}[c]
    \begin{column}{0.46\textwidth}
      \centering
      \begin{tikzpicture}[node distance=3mm]
        \node[caixa,minimum width=4.6cm] (d) {Seu código Dart};
        \node[caixapetroleo,minimum width=4.6cm,below=of d] (p) {Plugin (\cod{geolocator}, \cod{camera}\ldots)};
        \node[caixa,minimum width=2.2cm,below=5mm of p.south west,anchor=north west] (a) {Android\\ \tiny Kotlin/Java};
        \node[caixa,minimum width=2.2cm,below=5mm of p.south east,anchor=north east] (i) {iOS\\ \tiny Swift/ObjC};
        \node[caixaescura,minimum width=4.6cm,below=10mm of p] (h) {Hardware e sistema};
        \draw[seta] (d) -- (p);
        \draw[seta] (p.south -| a) -- node[rotulo,left]{\tiny canal} (a);
        \draw[seta] (p.south -| i) -- node[rotulo,right]{\tiny canal} (i);
        \draw[seta] (a) -- (a |- h.north); \draw[seta] (i) -- (i |- h.north);
      \end{tikzpicture}
    \end{column}
    \begin{column}{0.52\textwidth}
      \begin{conceito}[Plugin]
        Pacote com uma API em Dart e implementações nativas para cada plataforma, ligadas por \textbf{canais de plataforma} (Módulo 19 mostra como criar o seu).
      \end{conceito}
      \begin{dica}[Escolhendo um plugin no pub.dev]
        Prefira selo \textbf{Flutter Favorite} ou publicador verificado (\cod{flutter.dev}, \cod{fluttercommunity.dev}); confira plataformas suportadas, data da última versão, \emph{issues} abertas e licença.
      \end{dica}
      \begin{atencao}
        Recurso de hardware se testa em \textbf{aparelho real}. Emuladores simulam GPS e sensores, mas não Bluetooth, NFC ou USB.
      \end{atencao}
    \end{column}
  \end{columns}
\end{frame}

% =============================================================================
\section{Permissões}

\begin{frame}[fragile]{Declarar nas duas plataformas}
  \begin{columns}[T]
    \begin{column}{0.5\textwidth}
\begin{codigo}[title=android/app/src/main/AndroidManifest.xml]{XML}
<uses-permission
    android:name="android.permission.CAMERA" />
<uses-permission android:name=
    "android.permission.ACCESS_FINE_LOCATION" />
<uses-permission android:name=
    "android.permission.POST_NOTIFICATIONS" />
\end{codigo}
\begin{codigo}[title=ios/Runner/Info.plist]{XML}
<key>NSCameraUsageDescription</key>
<string>Fotografar equipamentos e ler QR codes.</string>
<key>NSLocationWhenInUseUsageDescription</key>
<string>Confirmar presença no laboratório.</string>
\end{codigo}
    \end{column}
    \begin{column}{0.48\textwidth}
      \begin{conceito}[Duas etapas]
        \textbf{Declarar}: manifesto (Android) e \cod{Info.plist} (iOS).\\
        \textbf{Pedir}: em tempo de execução, no momento em que o recurso é usado.
      \end{conceito}
      \begin{erro}[Rejeição na App Store]
        Falta de texto de uso no \cod{Info.plist} derruba o app no iOS; texto vago (``para funcionar'') reprova na revisão. Explique o benefício ao usuário.
      \end{erro}
    \end{column}
  \end{columns}
\end{frame}

\begin{frame}[fragile]{Pedindo com \cod{permission\_handler}}
  \begin{columns}[T]
    \begin{column}{0.58\textwidth}
\begin{codigo}[title=utils/permissoes.dart]{Dart}
Future<bool> garantir(Permission p,
    Future<bool> Function() explicar) async {
  var status = await p.status;
  if (status.isGranted || status.isLimited) return true;

  if (await p.shouldShowRequestRationale &&
      !await explicar()) return false;     // diálogo seu

  status = await p.request();
  if (status.isPermanentlyDenied) {
    await openAppSettings();               // só o usuário
    return false;
  }
  return status.isGranted || status.isLimited;
}

// várias de uma vez
final r = await [Permission.bluetoothScan,
    Permission.bluetoothConnect].request();
\end{codigo}
    \end{column}
    \begin{column}{0.4\textwidth}
      \begin{conceito}[Estados]
        \cod{granted}, \cod{denied}, \cod{permanentlyDenied}, \cod{limited} (acesso parcial, como só algumas fotos) e \cod{restricted} (controle parental).
      \end{conceito}
      \begin{atencao}[iOS]
        O \cod{permission\_handler} exige ativar cada permissão no \cod{ios/Podfile} (macros \cod{PERMISSION\_CAMERA=1} etc.).
      \end{atencao}
      \begin{dica}
        Peça no contexto: ao tocar em ``Ler QR code'', não na abertura do app.
      \end{dica}
    \end{column}
  \end{columns}
\end{frame}

% =============================================================================
\section{Câmera, galeria e QR code}

\begin{frame}[fragile]{Foto rápida com \cod{image\_picker}}
  \begin{columns}[T]
    \begin{column}{0.54\textwidth}
\begin{codigo}[title=Foto do equipamento]{Dart}
final _picker = ImagePicker();

Future<XFile?> fotografar() => _picker.pickImage(
      source: ImageSource.camera,
      maxWidth: 1600,          // reduz o arquivo
      imageQuality: 80,
    );

Future<List<XFile>> escolherDaGaleria() =>
    _picker.pickMultiImage(limit: 5);

// Android: o sistema pode encerrar o app
// enquanto a câmera está aberta
Future<void> recuperarPerdidos() async {
  final r = await _picker.retrieveLostData();
  if (!r.isEmpty) salvar(r.files!);
}
\end{codigo}
    \end{column}
    \begin{column}{0.44\textwidth}
      \begin{conceito}[Quando usar]
        Para tirar uma foto ou escolher imagens usando o app de câmera e o seletor do próprio sistema. Também grava vídeo com \cod{pickVideo}.
      \end{conceito}
      \begin{dica}[Galeria sem permissão]
        No Android, o \textbf{Photo Picker} do sistema dispensa permissão de armazenamento: o usuário escolhe e o app recebe só aquelas fotos.
      \end{dica}
      \begin{dica}
        \cod{XFile} vem do pacote \cod{cross\_file}: funciona também na web.
      \end{dica}
    \end{column}
  \end{columns}
\end{frame}

\begin{frame}[fragile]{Câmera dentro do app com \cod{camera}}
  \begin{columns}[T]
    \begin{column}{0.58\textwidth}
\begin{codigo}[title=Tela de captura (trechos do State)]{Dart}
CameraController? _ctrl;

Future<void> _iniciar() async {
  final cams = await availableCameras();
  final traseira = cams.firstWhere(
      (c) => c.lensDirection == CameraLensDirection.back);
  final ctrl = CameraController(traseira,
      ResolutionPreset.high, enableAudio: false);
  await ctrl.initialize();
  if (mounted) setState(() => _ctrl = ctrl);
}

@override
void didChangeAppLifecycleState(AppLifecycleState s) {
  if (s == AppLifecycleState.inactive) _ctrl?.dispose();
  if (s == AppLifecycleState.resumed) _iniciar();
}

// State com o mixin WidgetsBindingObserver
// build: CameraPreview(_ctrl!)
// botão: final foto = await _ctrl!.takePicture();
\end{codigo}
    \end{column}
    \begin{column}{0.4\textwidth}
      \begin{conceito}[Quando usar]
        Câmera como parte da interface: moldura sobre a imagem, fotos em sequência, quadros para visão computacional (\cod{startImageStream}).
      \end{conceito}
      \begin{atencao}[Ciclo de vida]
        Libere a câmera quando o app sai de cena e reabra ao voltar; outro app pode precisar dela.
      \end{atencao}
    \end{column}
  \end{columns}
\end{frame}

\begin{frame}[fragile]{QR code e código de barras}
  \begin{columns}[T]
    \begin{column}{0.56\textwidth}
\begin{codigo}[title=Lendo a etiqueta do equipamento]{Dart}
final _scanner = MobileScannerController(
    formats: const [BarcodeFormat.qrCode]);

@override
Widget build(BuildContext context) => MobileScanner(
  controller: _scanner,
  onDetect: (captura) {
    final valor = captura.barcodes.first.rawValue;
    if (valor == null) return;
    final uri = Uri.tryParse(valor);
    if (uri?.scheme == 'labmonitor') {
      _scanner.stop();
      context.go('/equipamentos/${uri!.pathSegments.last}');
    }
  },
);
\end{codigo}
\begin{codigo}[title=Gerando a etiqueta (qr\_flutter)]{Dart}
QrImageView(data: 'labmonitor://eq/42', size: 180)
\end{codigo}
    \end{column}
    \begin{column}{0.42\textwidth}
      \begin{conceito}[\cod{mobile\_scanner}]
        Lê QR code, EAN, Code 128 e outros formatos usando ML Kit (Android) e Vision (iOS), sem internet.
      \end{conceito}
      \begin{atencao}
        Valide o conteúdo lido: um QR code pode conter qualquer texto. Aceite só o seu esquema e IDs válidos.
      \end{atencao}
      \begin{dica}
        Pare o leitor depois da primeira leitura válida para não navegar várias vezes.
      \end{dica}
    \end{column}
  \end{columns}
\end{frame}

% =============================================================================
\section{Localização e mapas}

\begin{frame}[fragile]{GPS com \cod{geolocator}}
  \begin{columns}[T]
    \begin{column}{0.58\textwidth}
\begin{codigo}[title=data/services/localizacao\_service.dart]{Dart}
Future<Position> posicaoAtual() async {
  if (!await Geolocator.isLocationServiceEnabled()) {
    throw const LocalizacaoDesligada();
  }
  var p = await Geolocator.checkPermission();
  if (p == LocationPermission.denied) {
    p = await Geolocator.requestPermission();
  }
  if (p == LocationPermission.deniedForever) {
    throw const PermissaoNegada();   // openAppSettings()
  }
  return Geolocator.getCurrentPosition();
}

Stream<Position> acompanhar() =>
    Geolocator.getPositionStream(locationSettings:
        const LocationSettings(distanceFilter: 10));

bool noLab(Position p, Lab l) => Geolocator.distanceBetween(
    p.latitude, p.longitude, l.lat, l.lng) < 50;
\end{codigo}
    \end{column}
    \begin{column}{0.4\textwidth}
      \begin{conceito}[Precisão e bateria]
        \cod{LocationSettings} define a precisão (\cod{high} usa GPS e gasta mais) e o \cod{distanceFilter} (emitir só após andar N metros).
      \end{conceito}
      \begin{dica}[Check-in]
        \cod{noLab} é uma cerca virtual (\emph{geofence}) simples: presença só a menos de 50\,m.
      \end{dica}
      \begin{atencao}
        Localização em segundo plano exige permissão extra e justificativa nas lojas.
      \end{atencao}
    \end{column}
  \end{columns}
\end{frame}

\begin{frame}[fragile]{Mapas}
  \begin{columns}[T]
    \begin{column}{0.5\textwidth}
\begin{codigo}[title=google\_maps\_flutter]{Dart}
GoogleMap(
  initialCameraPosition: const CameraPosition(
      target: LatLng(-2.5560, -44.3080), zoom: 17),
  myLocationEnabled: true,
  markers: {
    const Marker(
      markerId: MarkerId('lab-3'),
      position: LatLng(-2.5560, -44.3080),
      infoWindow: InfoWindow(title: 'Lab. 3'),
    ),
  },
)
\end{codigo}
      \begin{atencao}
        Exige chave da Google Maps Platform, restrita ao seu app.
      \end{atencao}
    \end{column}
    \begin{column}{0.48\textwidth}
\begin{codigo}[title=flutter\_map + OpenStreetMap]{Dart}
FlutterMap(
  options: const MapOptions(
      initialCenter: LatLng(-2.5560, -44.3080),
      initialZoom: 17),
  children: [
    TileLayer(
      urlTemplate: 'https://tile.openstreetmap'
          '.org/{z}/{x}/{y}.png',
      userAgentPackageName: 'br.edu.labmonitor'),
    const MarkerLayer(markers: [Marker(
      point: LatLng(-2.5560, -44.3080),
      child: Icon(Icons.science))]),
  ],
)
\end{codigo}
      \begin{dica}
        Sem chave; mostre a atribuição do OSM e, com muito tráfego, use outro provedor.
      \end{dica}
    \end{column}
  \end{columns}
\end{frame}

% =============================================================================
\section{Sensores e dados do aparelho}

\begin{frame}[fragile]{Sensores de movimento com \cod{sensors\_plus}}
  \begin{columns}[T]
    \begin{column}{0.56\textwidth}
\begin{codigo}[title=Agitar para atualizar]{Dart}
StreamSubscription<UserAccelerometerEvent>? _sub;
DateTime _ultimo = DateTime(0);

void detectarAgitar(VoidCallback aoAgitar) {
  _sub = userAccelerometerEventStream(
    samplingPeriod: SensorInterval.gameInterval,
  ).listen((e) {
    final a = sqrt(e.x * e.x + e.y * e.y + e.z * e.z);
    final agora = DateTime.now();
    if (a > 15 && agora.difference(_ultimo).inSeconds > 1) {
      _ultimo = agora;
      aoAgitar();
    }
  });
}

void parar() => _sub?.cancel();
\end{codigo}
    \end{column}
    \begin{column}{0.42\textwidth}
      \scriptsize
      \begin{tabularx}{\linewidth}{@{}l L@{}}
        \toprule
        \cabecalho{Stream} & \cabecalho{Mede} \\
        \midrule
        \cod{accelerometer\ldots} & aceleração com gravidade (m/s\textsuperscript{2}) \\
        \cod{userAccelerometer\ldots} & aceleração sem gravidade \\
        \cod{gyroscope\ldots} & rotação (rad/s) \\
        \cod{magnetometer\ldots} & campo magnético ($\mu$T) \\
        \cod{barometer\ldots} & pressão (hPa) \\
        \bottomrule
      \end{tabularx}
      \begin{atencao}
        Sensores gastam bateria: cancele a assinatura quando a tela sai e use a menor frequência que resolve.
      \end{atencao}
    \end{column}
  \end{columns}
\end{frame}

\begin{frame}[fragile]{Bateria, rede e dados do aparelho}
  \begin{columns}[T]
    \begin{column}{0.56\textwidth}
\begin{codigo}[title=Pacotes ``plus'' da comunidade]{Dart}
// battery_plus
final bateria = Battery();
final nivel = await bateria.batteryLevel;          // 0..100
bateria.onBatteryStateChanged.listen((s) =>
    carregando = s == BatteryState.charging);

// connectivity_plus
Connectivity().onConnectivityChanged.listen((tipos) =>
    offline = tipos.contains(ConnectivityResult.none));

// device_info_plus e package_info_plus
final info = await DeviceInfoPlugin().androidInfo;
log('${info.model} - Android SDK ${info.version.sdkInt}');
final app = await PackageInfo.fromPlatform();
log('LabMonitor ${app.version}+${app.buildNumber}');

// wakelock_plus: tela sempre ligada no modo painel
WakelockPlus.enable();
\end{codigo}
    \end{column}
    \begin{column}{0.42\textwidth}
      \begin{conceito}[Usos no LabMonitor]
        \begin{itemize}
          \item Reduzir a frequência de leitura com bateria baixa.
          \item Mostrar a faixa ``sem conexão''.
          \item Enviar modelo, sistema e versão do app nos relatórios de erro.
          \item Manter o tablet do laboratório com a tela acesa.
        \end{itemize}
      \end{conceito}
      \begin{dica}
        Estes pacotes são mantidos pela \textbf{Flutter Community} (\cod{plus\_plugins}).
      \end{dica}
    \end{column}
  \end{columns}
\end{frame}

% =============================================================================
\section{Biometria e NFC}

\begin{frame}[fragile]{Biometria com \cod{local\_auth}}
  \begin{columns}[T]
    \begin{column}{0.56\textwidth}
\begin{codigo}[title=Confirmar uma ação sensível (local\_auth 3)]{Dart}
final _auth = LocalAuthentication();

Future<bool> confirmarIdentidade() async {
  final pode = await _auth.canCheckBiometrics ||
      await _auth.isDeviceSupported();
  if (!pode) return false;
  try {
    return await _auth.authenticate(
      localizedReason: 'Confirme para liberar a bancada',
      biometricOnly: false,       // aceita PIN/padrão
      persistAcrossBackgrounding: true,
    );
  } on LocalAuthException {
    return false;                 // cancelado, bloqueado...
  }
}
\end{codigo}
    \end{column}
    \begin{column}{0.42\textwidth}
      \begin{atencao}[Configuração]
        Android: a \cod{MainActivity} deve estender \cod{FlutterFragmentActivity}. iOS: \cod{NSFaceIDUsageDescription} no \cod{Info.plist}.
      \end{atencao}
      \begin{conceito}[O que a biometria prova]
        Só que o dono do aparelho está presente. Ela \textbf{desbloqueia} um segredo local (Módulo 10); não substitui o login no servidor.
      \end{conceito}
    \end{column}
  \end{columns}
\end{frame}

\begin{frame}[fragile]{NFC: etiquetas nas bancadas}
  \begin{columns}[T]
    \begin{column}{0.56\textwidth}
\begin{codigo}[title=Check-in por aproximação (flutter\_nfc\_kit)]{Dart}
Future<String?> lerEtiqueta() async {
  final disp = await FlutterNfcKit.nfcAvailability;
  if (disp != NFCAvailability.available) return null;
  try {
    final tag = await FlutterNfcKit.poll(
      timeout: const Duration(seconds: 15),
      iosAlertMessage: 'Aproxime da etiqueta da bancada',
    );
    return tag.id;            // identificador único da tag
  } finally {
    await FlutterNfcKit.finish(iosAlertMessage: 'Pronto!');
  }
}
\end{codigo}
      \begin{dica}
        Para ler ou gravar dados na tag (uma URL, um texto), use \cod{readNDEFRecords} e \cod{writeNDEFRecords}.
      \end{dica}
    \end{column}
    \begin{column}{0.42\textwidth}
      \begin{conceito}[NFC]
        Comunicação a poucos centímetros. Etiquetas NTAG custam centavos e não precisam de bateria: ótimas para identificar bancadas e equipamentos.
      \end{conceito}
      \begin{atencao}[Configuração]
        Android: permissão \cod{android.permission.NFC}. iOS: capacidade \emph{NFC Tag Reading} no Xcode e \cod{NFCReaderUsageDescription}; o iOS mostra uma folha de leitura própria.
      \end{atencao}
    \end{column}
  \end{columns}
\end{frame}

% =============================================================================
\section{Notificações e segundo plano}

\begin{frame}[fragile]{Notificações locais}
  \begin{columns}[T]
    \begin{column}{0.58\textwidth}
\begin{codigo}[title=flutter\_local\_notifications (versão 20+)]{Dart}
final _notif = FlutterLocalNotificationsPlugin();
const _android = AndroidInitializationSettings('@mipmap/ic_launcher');
Future<void> iniciar() async {
  await _notif.initialize(
    settings: const InitializationSettings(android: _android,
        iOS: DarwinInitializationSettings()),
    onDidReceiveNotificationResponse: (r) =>
        router.go(r.payload ?? '/'));
  await _notif.resolvePlatformSpecificImplementation<
          AndroidFlutterLocalNotificationsPlugin>()
      ?.requestNotificationsPermission();   // Android 13+
}

const _alertas = NotificationDetails(
  android: AndroidNotificationDetails('alertas', 'Alertas',
      importance: Importance.high, priority: Priority.high));
Future<void> alertar(int b, double t) => _notif.show(
    id: b, title: 'Bancada $b', body: 'Temperatura: $t °C',
    notificationDetails: _alertas, payload: '/bancadas/$b');
\end{codigo}
    \end{column}
    \begin{column}{0.4\textwidth}
      \begin{conceito}[Canais (Android 8+)]
        Cada notificação pertence a um canal (\cod{'alertas'}); o usuário pode silenciar um canal sem perder os outros.
      \end{conceito}
      \begin{dica}
        \cod{zonedSchedule} agenda lembretes (pacote \cod{timezone}).
      \end{dica}
      \begin{atencao}
        Desde a versão 20 os parâmetros são nomeados; exemplos antigos não compilam.
      \end{atencao}
    \end{column}
  \end{columns}
\end{frame}

\begin{frame}[fragile]{Trabalho em segundo plano}
  \begin{columns}[T]
    \begin{column}{0.54\textwidth}
\begin{codigo}[title=Sincronização periódica (workmanager)]{Dart}
@pragma('vm:entry-point')
void despachante() {
  Workmanager().executeTask((tarefa, dados) async {
    WidgetsFlutterBinding.ensureInitialized();
    await sincronizarLeituras();        // Módulo 11
    return true;                        // sucesso
  });
}

Future<void> agendar() async {
  await Workmanager().initialize(despachante);
  await Workmanager().registerPeriodicTask(
    'sync-leituras', 'sincronizar',
    frequency: const Duration(minutes: 15),
    constraints: Constraints(
        networkType: NetworkType.connected),
  );
}
\end{codigo}
    \end{column}
    \begin{column}{0.44\textwidth}
      \scriptsize
      \begin{tabularx}{\linewidth}{@{}l L@{}}
        \toprule
        \cabecalho{Necessidade} & \cabecalho{Solução} \\
        \midrule
        tarefa adiável & WorkManager (Android, mín. 15\,min); BGTaskScheduler (iOS, quando o sistema quiser) \\
        tarefa contínua & \emph{foreground service} com notificação fixa (\cod{flutter\_foreground\_task}) \\
        BLE em segundo plano & iOS: modo \cod{bluetooth-central} \\
        aviso do servidor & push com FCM (Módulo 12) \\
        \bottomrule
      \end{tabularx}
      \begin{atencao}
        O sistema decide quando rodar e encerra apps que gastam bateria. No Android 14+, todo \emph{foreground service} declara um tipo (\cod{connectedDevice}, \cod{location}\ldots).
      \end{atencao}
    \end{column}
  \end{columns}
\end{frame}

% =============================================================================
\section{Arquivos, links e áudio}

\begin{frame}[fragile]{Compartilhar, escolher arquivos e abrir links}
  \begin{columns}[T]
    \begin{column}{0.56\textwidth}
\begin{codigo}[title={share\_plus, file\_picker e url\_launcher}]{Dart}
// enviar o relatório por qualquer app instalado
await SharePlus.instance.share(ShareParams(
  text: 'Relatório do Lab. 3 - outubro',
  files: [XFile(caminhoCsv)],
));

// importar uma planilha de equipamentos
final r = await FilePicker.platform.pickFiles(
    type: FileType.custom, allowedExtensions: ['csv']);
final arquivo = r?.files.single.path;

// telefone, e-mail, mapa e navegador
await launchUrl(Uri.parse('tel:+559832180000'));
await launchUrl(Uri.parse('mailto:lab3@ifma.edu.br'));
await launchUrl(Uri.parse('https://labmonitor.dev'),
    mode: LaunchMode.externalApplication);
\end{codigo}
    \end{column}
    \begin{column}{0.42\textwidth}
      \begin{conceito}[\emph{Deep links}]
        O \cod{go\_router} abre a rota certa quando o app é aberto por um link. Com \textbf{App Links} (Android) e \textbf{Universal Links} (iOS), \cod{https://labmonitor.dev/eq/42} abre o app em vez do navegador.
      \end{conceito}
      \begin{atencao}
        Android 11+: declare em \cod{<queries>} os esquemas (\cod{tel}, \cod{mailto}, \cod{https}) usados com \cod{canLaunchUrl}.
      \end{atencao}
      \begin{dica}
        A API antiga \cod{Share.shareXFiles} foi substituída por \cod{SharePlus.instance.share}.
      \end{dica}
    \end{column}
  \end{columns}
\end{frame}

\begin{frame}[fragile]{Vibração, som, voz e microfone}
  \begin{columns}[T]
    \begin{column}{0.56\textwidth}
\begin{codigo}[title=Alarme de temperatura]{Dart}
HapticFeedback.heavyImpact();        // já vem no Flutter

final player = AudioPlayer();        // just_audio
await player.setAsset('assets/sons/alarme.mp3');
await player.setLoopMode(LoopMode.one);
player.play();

final tts = FlutterTts();            // flutter_tts
await tts.setLanguage('pt-BR');
await tts.speak('Temperatura alta na bancada três');

final gravador = AudioRecorder();    // record
if (await gravador.hasPermission()) {
  await gravador.start(const RecordConfig(), path: destino);
}
final arquivo = await gravador.stop();
\end{codigo}
    \end{column}
    \begin{column}{0.42\textwidth}
      \scriptsize
      \begin{tabularx}{\linewidth}{@{}l L@{}}
        \toprule
        \cabecalho{Pacote} & \cabecalho{Para quê} \\
        \midrule
        \cod{HapticFeedback} & toque tátil curto \\
        \cod{vibration} & padrões longos de vibração \\
        \cod{just\_audio} & tocar áudio local ou remoto \\
        \cod{flutter\_tts} & texto em fala \\
        \cod{speech\_to\_text} & fala em texto \\
        \cod{record} & gravar do microfone \\
        \bottomrule
      \end{tabularx}
      \begin{atencao}
        O microfone exige permissão nas duas plataformas (\cod{RECORD\_AUDIO} no Android).
      \end{atencao}
      \begin{dica}
        Alertas importantes: som, vibração \textbf{e} texto.
      \end{dica}
    \end{column}
  \end{columns}
\end{frame}

% =============================================================================
\section{Bluetooth Low Energy com ESP32}

\begin{frame}{Conceitos de BLE}
  \begin{columns}[c]
    \begin{column}{0.5\textwidth}
      \centering
      \begin{tikzpicture}[node distance=3mm]
        \node[caixa,minimum width=2.4cm] (app) {App\\ \tiny \textbf{central}};
        \node[caixaescura,minimum width=2.6cm,right=12mm of app] (dev) {ESP32\\ \tiny \textbf{periférico}};
        \draw[setaambar,<->] (app) -- node[rotulo,above]{\tiny GATT} (dev);
        \node[caixapetroleo,minimum width=5cm,below=6mm of $(app)!0.5!(dev)$] (svc) {\textbf{Serviço} \tiny UUID 4fafc201-\ldots{} (Ambiente)};
        \node[caixaambar,minimum width=4.4cm,below=of svc] (c1) {\textbf{Característica} Temperatura\\ \tiny UUID beb5483e-\ldots{} \quad READ $\cdot$ NOTIFY};
        \node[caixaambar,minimum width=4.4cm,below=of c1] (c2) {\textbf{Característica} Comando\\ \tiny WRITE};
      \end{tikzpicture}
    \end{column}
    \begin{column}{0.48\textwidth}
      \begin{itemize}
        \item \termo{Central} procura e conecta; \termo{periférico} anuncia (\emph{advertising}) e serve dados.
        \item \termo{GATT}: dados organizados em \textbf{serviços} e \textbf{características}, cada um com um UUID.
        \item Operações: \textbf{read}, \textbf{write} e \textbf{notify} (o periférico envia quando o valor muda).
        \item \termo{MTU}: tamanho máximo de cada pacote (23 a 512 bytes).
      \end{itemize}
      \begin{atencao}
        BLE é diferente do Bluetooth clássico: módulos HC-05/HC-06 do Arduino \textbf{não} são BLE. Use ESP32, HM-10 ou nRF52.
      \end{atencao}
    \end{column}
  \end{columns}
\end{frame}

\begin{frame}[fragile]{Firmware do ESP32 (Arduino)}
  \begin{columns}[T]
    \begin{column}{0.56\textwidth}
\begin{codigo}[title=labmonitor\_ble.ino --- configuração]{CPP}
#include <BLEDevice.h>
#include <BLE2902.h>
#define SVC_UUID  "4fafc201-1fb5-459e-8fcc-c5c9c331914b"
#define TEMP_UUID "beb5483e-36e1-4688-b7f5-ea07361b26a8"
BLECharacteristic *temp;
void setup() {
  BLEDevice::init("LabMonitor-B3");
  BLEServer *srv = BLEDevice::createServer();
  BLEService *svc = srv->createService(SVC_UUID);
  temp = svc->createCharacteristic(TEMP_UUID,
    BLECharacteristic::PROPERTY_READ |
    BLECharacteristic::PROPERTY_NOTIFY);
  temp->addDescriptor(new BLE2902());   // CCCD
  svc->start();
  BLEDevice::getAdvertising()->addServiceUUID(SVC_UUID);
  BLEDevice::startAdvertising();
}
\end{codigo}
    \end{column}
    \begin{column}{0.42\textwidth}
\begin{codigo}[title=--- laço principal]{CPP}
void loop() {
  // centésimos de grau: 24,87 °C -> 2487
  int16_t c = round(lerTemperatura() * 100);
  temp->setValue((uint8_t*)&c, sizeof(c));
  temp->notify();
  delay(1000);
}
\end{codigo}
      \begin{conceito}[Formato binário]
        \cod{int16\_t} em centésimos de grau: 2 bytes, \emph{little-endian}.
      \end{conceito}
      \begin{atencao}[Contrato]
        Documente UUIDs, escala e ordem dos bytes; gere UUIDs próprios.
      \end{atencao}
    \end{column}
  \end{columns}
\end{frame}

\begin{frame}[fragile]{Permissões e licença do \cod{flutter\_blue\_plus}}
  \begin{columns}[T]
    \begin{column}{0.56\textwidth}
\begin{codigo}[title=AndroidManifest.xml]{XML}
<uses-permission android:name=
    "android.permission.BLUETOOTH_SCAN"
    android:usesPermissionFlags="neverForLocation" />
<uses-permission android:name=
    "android.permission.BLUETOOTH_CONNECT" />
<!-- Android 11 ou anterior -->
<uses-permission android:name=
    "android.permission.ACCESS_FINE_LOCATION"
    android:maxSdkVersion="30" />
\end{codigo}
\begin{codigo}[title=Info.plist]{XML}
<key>NSBluetoothAlwaysUsageDescription</key>
<string>Conectar aos sensores das bancadas.</string>
\end{codigo}
    \end{column}
    \begin{column}{0.42\textwidth}
      \begin{atencao}[Licença do pacote]
        Desde a versão 2.0 (2025) o \cod{flutter\_blue\_plus} usa licença própria, que já mudou mais de uma vez: uso pessoal, educacional e sem fins lucrativos é livre; uso comercial pode exigir licença paga, e o pacote envia telemetria de licença no \emph{build}. Leia o \cod{LICENSE} da versão que adotar.
      \end{atencao}
      \begin{dica}[Alternativas]
        \cod{universal\_ble} e \cod{bluetooth\_low\_energy}; compare licença, plataformas e manutenção.
      \end{dica}
    \end{column}
  \end{columns}
\end{frame}

\begin{frame}[fragile]{Procurando o ESP32}
  \begin{columns}[T]
    \begin{column}{0.58\textwidth}
\begin{codigo}[title=data/services/ble\_service.dart]{Dart}
final svcUuid = Guid('4fafc201-1fb5-459e-8fcc-c5c9c331914b');
final tempUuid = Guid('beb5483e-36e1-4688-b7f5-ea07361b26a8');

Future<List<ScanResult>> procurar() async {
  await FlutterBluePlus.adapterState
      .where((s) => s == BluetoothAdapterState.on).first;
  final achados = <DeviceIdentifier, ScanResult>{};
  final sub = FlutterBluePlus.onScanResults.listen((rs) {
    for (final r in rs) achados[r.device.remoteId] = r;
  });
  FlutterBluePlus.cancelWhenScanComplete(sub);
  await FlutterBluePlus.startScan(
      withServices: [svcUuid],
      timeout: const Duration(seconds: 10));
  await FlutterBluePlus.isScanning
      .where((v) => v == false).first;
  return achados.values.toList();
}
\end{codigo}
    \end{column}
    \begin{column}{0.4\textwidth}
      \begin{conceito}[\cod{ScanResult}]
        Traz o dispositivo, o nome anunciado (\cod{advertisementData.advName}) e o \textbf{RSSI} (força do sinal, que indica a proximidade).
      \end{conceito}
      \begin{dica}
        O filtro \cod{withServices} faz o sistema entregar só os seus dispositivos e economiza bateria.
      \end{dica}
      \begin{atencao}
        Peça \cod{bluetoothScan} e \cod{bluetoothConnect} antes de procurar e trate o Bluetooth desligado.
      \end{atencao}
    \end{column}
  \end{columns}
\end{frame}

\begin{frame}[fragile]{Conectando e recebendo a temperatura}
  \begin{columns}[T]
    \begin{column}{0.58\textwidth}
\begin{codigo}[title=Notify vira Stream<double>]{Dart}
Stream<double> temperatura(BluetoothDevice d) async* {
  await d.connect();
  final servicos = await d.discoverServices();
  final c = servicos
      .firstWhere((s) => s.uuid == svcUuid)
      .characteristics
      .firstWhere((c) => c.uuid == tempUuid);
  await c.setNotifyValue(true);
  yield* c.onValueReceived.map((bytes) {
    final v = ByteData.sublistView(Uint8List.fromList(bytes))
        .getInt16(0, Endian.little);
    return v / 100;                 // centésimos -> °C
  });
}

// enviar um comando
await comando.write(utf8.encode('LED=1'));
// ao sair: await d.disconnect();
\end{codigo}
    \end{column}
    \begin{column}{0.4\textwidth}
      \begin{conceito}[Fluxo]
        Conectar $\rightarrow$ descobrir serviços $\rightarrow$ achar a característica $\rightarrow$ ativar notify $\rightarrow$ converter os bytes.
      \end{conceito}
      \begin{atencao}[Desconexões]
        Observe \cod{d.connectionState}; ao reconectar, descubra os serviços de novo. Use \cod{d.cancelWhenDisconnected(sub)} para não acumular assinaturas.
      \end{atencao}
      \begin{dica}
        O \cod{Stream} entra no view model como qualquer outro (Módulos 06 e 08).
      \end{dica}
    \end{column}
  \end{columns}
\end{frame}

% =============================================================================
\section{Serial USB com Arduino}

\begin{frame}[fragile]{Arduino pelo cabo USB (Android)}
  \begin{columns}[T]
    \begin{column}{0.4\textwidth}
\begin{codigo}[title=sensor\_serial.ino]{CPP}
void setup() {
  Serial.begin(115200);
  pinMode(13, OUTPUT);
}
void loop() {
  float t = analogRead(A0)
      * (5.0 / 1023.0) * 100;  // LM35
  Serial.print("T=");
  Serial.println(t, 2);
  if (Serial.available()) {
    String cmd =
        Serial.readStringUntil('\n');
    digitalWrite(13,
        cmd == "LED=1" ? HIGH : LOW);
  }
  delay(1000);
}
\end{codigo}
    \end{column}
    \begin{column}{0.58\textwidth}
\begin{codigo}[title=Lendo linhas com usb\_serial]{Dart}
Future<Stream<String>> abrirSerial() async {
  final dispositivos = await UsbSerial.listDevices();
  final porta = await dispositivos.first.create();
  if (porta == null || !await porta.open()) {
    throw const SerialIndisponivel();
  }
  await porta.setDTR(true);
  await porta.setRTS(true);
  await porta.setPortParameters(115200, UsbPort.DATABITS_8,
      UsbPort.STOPBITS_1, UsbPort.PARITY_NONE);
  return porta.inputStream!
      .cast<List<int>>()
      .transform(utf8.decoder)
      .transform(const LineSplitter());      // "T=24.87"
}
\end{codigo}
      \begin{atencao}
        Só no Android, com cabo \textbf{OTG}; o sistema pede autorização ao conectar. O iOS não permite serial USB genérica.
      \end{atencao}
    \end{column}
  \end{columns}
\end{frame}

% =============================================================================
\section{MQTT e escolha do protocolo}

\begin{frame}{MQTT: publicar e assinar}
  \begin{center}
  \begin{tikzpicture}[node distance=8mm]
    \node[caixaescura,minimum width=2.6cm] (e1) {ESP32 bancada 3\\ \tiny publica};
    \node[caixaescura,minimum width=2.6cm,below=4mm of e1] (e2) {ESP32 bancada 4\\ \tiny publica};
    \node[caixaambarforte,minimum width=2.8cm,minimum height=14mm,right=16mm of $(e1)!0.5!(e2)$] (b) {\textbf{Broker}\\ \tiny Mosquitto, EMQX, HiveMQ};
    \node[caixa,minimum width=2.6cm,right=16mm of b,yshift=6mm] (app) {App LabMonitor\\ \tiny assina \cod{ifma/dee/+/+/temp}};
    \node[caixa,minimum width=2.6cm,right=16mm of b,yshift=-6mm] (db) {Servidor / histórico\\ \tiny assina \cod{ifma/\#}};
    \draw[seta] (e1) -- node[rotulo,above,sloped]{\cod{ifma/dee/lab1/b3/temp}} (b);
    \draw[seta] (e2) -- (b);
    \draw[seta] (b) -- (app); \draw[seta] (b) -- (db);
  \end{tikzpicture}
  \end{center}
  \begin{columns}[T]
    \begin{column}{0.32\textwidth}
      \begin{conceito}[Tópicos]
        Hierarquia com \cod{/}. Curingas: \cod{+} (um nível) e \cod{\#} (todos os níveis abaixo).
      \end{conceito}
    \end{column}
    \begin{column}{0.32\textwidth}
      \begin{conceito}[QoS]
        \textbf{0}: no máximo uma vez. \textbf{1}: pelo menos uma vez. \textbf{2}: exatamente uma vez.
      \end{conceito}
    \end{column}
    \begin{column}{0.32\textwidth}
      \begin{conceito}[Retained / LWT]
        \emph{Retained} guarda o último valor do tópico. \emph{Last Will} avisa quando a placa cai.
      \end{conceito}
    \end{column}
  \end{columns}
\end{frame}

\begin{frame}[fragile]{MQTT no app}
  \begin{columns}[T]
    \begin{column}{0.58\textwidth}
\begin{codigo}[title=data/services/mqtt\_service.dart (mqtt\_client)]{Dart}
Future<Stream<(String, double)>> leituras(String f) async {
  final c = MqttServerClient.withPort(
      'broker.labmonitor.ifma.edu.br', clienteId, 8883)
    ..secure = true                     // TLS
    ..keepAlivePeriod = 30
    ..autoReconnect = true;
  await c.connect(usuario, senha);
  c.subscribe(f, MqttQos.atLeastOnce);
  return c.updates!.expand((lote) => lote).map((m) {
    final pub = m.payload as MqttPublishMessage;
    final txt = MqttPublishPayload
        .bytesToStringAsString(pub.payload.message);
    return (m.topic, double.parse(txt));
  });
}
\end{codigo}
    \end{column}
    \begin{column}{0.4\textwidth}
\begin{codigo}[title=ESP32 (PubSubClient)]{CPP}
mqtt.setServer(BROKER, 1883);
mqtt.connect("esp32-b3", USER, PASS,
  "ifma/dee/lab1/b3/status", 1, true,
  "offline");                 // LWT
mqtt.publish("ifma/dee/lab1/b3/temp",
  String(t, 2).c_str(), true); // retained
\end{codigo}
      \begin{atencao}[Segurança]
        Broker público só para testes. Em produção: TLS, autenticação e regras de acesso por tópico.
      \end{atencao}
      \begin{dica}
        Publicar: \cod{publishMessage}.
      \end{dica}
    \end{column}
  \end{columns}
\end{frame}

\begin{frame}{Qual protocolo usar?}
  \footnotesize
  \begin{tabularx}{\textwidth}{@{}l L L L L@{}}
    \toprule
    & \cabecalho{BLE} & \cabecalho{USB serial} & \cabecalho{HTTP na rede local} & \cabecalho{MQTT via broker} \\
    \midrule
    Alcance & $\sim$10--30\,m & o cabo & a rede Wi-Fi & qualquer lugar com internet \\
    Consumo no sensor & muito baixo & alimentado pelo cabo & alto (Wi-Fi) & médio \\
    Tempo real & sim (notify) & sim & não (\emph{polling}) & sim (push) \\
    Plataformas & Android, iOS & só Android & todas & todas \\
    Vários clientes & um central por vez & um & sim & sim, muitos \\
    Caso ideal & configurar a placa, bancada ao lado & bancada fixa, Arduino sem rádio & painel na rede do laboratório & monitoramento contínuo e alertas \\
    \bottomrule
  \end{tabularx}
  \vspace{4pt}
  \begin{dica}[Arquitetura combinada]
    Use BLE para o \emph{provisionamento} (enviar SSID e senha do Wi-Fi ao ESP32 pelo app) e MQTT para a telemetria do dia a dia.
  \end{dica}
\end{frame}

% =============================================================================
\section{Arquitetura e referência rápida}

\begin{frame}[fragile]{Hardware atrás de uma interface}
  \begin{columns}[T]
    \begin{column}{0.54\textwidth}
\begin{codigo}[title=data/services/fonte\_leituras.dart]{Dart}
abstract interface class FonteLeituras {
  Stream<Leitura> leituras(String bancada);
}

class FonteBle implements FonteLeituras { ... }
class FonteMqtt implements FonteLeituras { ... }
class FonteSerial implements FonteLeituras { ... }

class FonteFake implements FonteLeituras {
  @override
  Stream<Leitura> leituras(String b) => Stream.periodic(
      const Duration(seconds: 1),
      (i) => Leitura(b, 22 + sin(i / 5) * 3));
}
\end{codigo}
    \end{column}
    \begin{column}{0.44\textwidth}
      \centering
      \begin{tikzpicture}[node distance=3mm]
        \node[caixa,minimum width=3.6cm] (vm) {View model};
        \node[caixapetroleo,minimum width=3.6cm,below=of vm] (r) {LeituraRepository};
        \node[caixaambar,minimum width=3.6cm,below=of r] (f) {\cod{FonteLeituras}};
        \node[caixa,font=\tiny,below=5mm of f,xshift=-13mm] (a) {BLE};
        \node[caixa,font=\tiny,below=5mm of f] (b) {MQTT};
        \node[caixa,font=\tiny,below=5mm of f,xshift=13mm] (c) {Fake};
        \draw[seta] (vm) -- (r); \draw[seta] (r) -- (f);
        \draw[seta] (f) -- (a); \draw[seta] (f) -- (b); \draw[seta] (f) -- (c);
      \end{tikzpicture}
      \begin{conceito}[Por quê]
        A UI e as regras não sabem de onde vem o dado. O \emph{fake} permite desenvolver sem placa, demonstrar em sala e testar (Módulo 14).
      \end{conceito}
    \end{column}
  \end{columns}
\end{frame}

\begin{frame}{Referência rápida: pacotes e permissões}
  \scriptsize
  \begin{tabularx}{\textwidth}{@{}p{2.0cm} p{3.0cm} p{3.4cm} L@{}}
    \toprule
    \cabecalho{Recurso} & \cabecalho{Pacote} & \cabecalho{Android (manifesto)} & \cabecalho{iOS (\cod{Info.plist})} \\
    \midrule
    Câmera e QR code & \cod{camera}, \cod{mobile\_scanner} & \cod{CAMERA} & \cod{NSCameraUsageDescription} \\
    Galeria & \cod{image\_picker} & nenhuma (Photo Picker) & \cod{NSPhotoLibraryUsageDescription} \\
    Localização & \cod{geolocator} & \cod{ACCESS\_FINE\_LOCATION} & \cod{NSLocationWhenInUseUsageDescription} \\
    Biometria & \cod{local\_auth} & \cod{USE\_BIOMETRIC} & \cod{NSFaceIDUsageDescription} \\
    NFC & \cod{flutter\_nfc\_kit} & \cod{NFC} & \cod{NFCReaderUsageDescription} + capacidade \\
    Notificações & \cod{flutter\_local\_}\-\cod{notifications} & \cod{POST\_NOTIFICATIONS} & pedido em tempo de execução \\
    Bluetooth LE & \cod{flutter\_blue\_plus} & \cod{BLUETOOTH\_SCAN} e \cod{\_CONNECT} & \cod{NSBluetoothAlwaysUsageDescription} \\
    Microfone & \cod{record} & \cod{RECORD\_AUDIO} & \cod{NSMicrophoneUsageDescription} \\
    Rede e MQTT & \cod{dio}, \cod{mqtt\_client} & \cod{INTERNET} & --- \\
    Serial USB & \cod{usb\_serial} & diálogo do sistema & não suportado \\
    \bottomrule
  \end{tabularx}
\end{frame}

% =============================================================================
\section{Fechamento}

\begin{frame}{Resumo}
  \begin{resumo}
    \begin{itemize}
      \item Plugins ligam Dart ao nativo; declare permissões e peça no contexto certo.
      \item Câmera, QR code, GPS, mapas e sensores viram \cod{Future}s e \cod{Stream}s em Dart.
      \item Biometria desbloqueia segredos locais; NFC identifica objetos por aproximação.
      \item Notificações locais, WorkManager e \emph{foreground services} cobrem o segundo plano.
      \item ESP32 e Arduino conversam por BLE, USB serial ou MQTT, atrás de uma interface comum.
    \end{itemize}
  \end{resumo}
\end{frame}

\begin{frame}{Glossário (1/2)}
  \begin{glossario}
    \emph{Plugin} & pacote com API em Dart e implementação nativa em cada plataforma \\
    Permissão em tempo de execução & autorização dada pelo usuário num diálogo do sistema \\
    ADC & conversor que transforma uma tensão em número \\
    Taxa de amostragem & quantas medidas são feitas por segundo \\
    Sensor MEMS & sensor micromecânico gravado no silício (acelerômetro, giroscópio) \\
    Fusão de sensores & combinação de sensores para compensar os defeitos de cada um \\
    GNSS & sistemas de navegação por satélite: GPS, GLONASS, Galileo, BeiDou \\
    Trilateração & cálculo da posição a partir das distâncias a pontos conhecidos \\
    Doze & modo que adia rede e alarmes com o aparelho parado \\
  \end{glossario}
\end{frame}

\begin{frame}{Glossário (2/2)}
  \begin{glossario}
    BLE & Bluetooth de baixo consumo, feito para sensores \\
    GATT & organização dos dados BLE em serviços e características \\
    Característica & valor com UUID que pode ser lido, escrito ou notificado \\
    MTU & tamanho máximo de um pacote ATT (23 bytes por padrão) \\
    MQTT & protocolo leve de publicação e assinatura, com \emph{broker} \\
    QoS & nível de garantia de entrega do MQTT: 0, 1 ou 2 \\
    UART & comunicação serial assíncrona (formato 8N1, \emph{baud rate}) \\
    NDEF & formato padrão das mensagens em etiquetas NFC \\
  \end{glossario}
\end{frame}

\begin{frame}{Questões de revisão}
  \begin{questoes}
    \item Descreva o caminho de uma temperatura do sensor do ESP32 até a tela. Onde entram o ADC e a taxa de amostragem?
    \item Por que o acelerômetro sozinho não basta para saber a orientação do celular?
    \item Quantos satélites o GPS precisa e por quê? Por que, dentro do laboratório, a posição piora?
    \item Quando pedir uma permissão, e o que fazer se o usuário negar?
    \item Por que o BLE consome tão pouco? O que é o intervalo de conexão?
    \item Que QoS do MQTT usar para leituras a cada segundo e para o comando que desliga o relé? Por quê?
    \item Explique o formato 8N1 e o que acontece se os dois lados usarem \emph{baud rates} diferentes.
    \item Como reduzir o consumo de bateria de um app que monitora a bancada o dia todo?
  \end{questoes}
\end{frame}

\begin{frame}{Exercícios}
  \begin{exercicio}[Exercício 1 --- Inventário]
    Leia o QR code de um equipamento, tire uma foto, registre a posição GPS e compartilhe um resumo.
  \end{exercicio}
  \begin{exercicio}[Exercício 2 --- Check-in]
    Confirme presença no laboratório exigindo distância menor que 50\,m \textbf{ou} leitura da etiqueta NFC da bancada, e biometria antes de liberar.
  \end{exercicio}
  \begin{exercicio}[Exercício 3 --- IoT]
    Grave o firmware no ESP32, receba a temperatura por BLE e dispare uma notificação local acima de 30\,°C.
  \end{exercicio}
  \begin{projeto}[LabMonitor --- Etapa 13]
    \cod{FonteLeituras} com implementações BLE, MQTT e \emph{fake}; check-in por QR/NFC; alarme sonoro e notificação de temperatura.
  \end{projeto}
\end{frame}

\begin{frame}{Para aprofundar}
  \small
  \begin{itemize}
    \item Plugins e canais de plataforma: \url{https://docs.flutter.dev/platform-integration/platform-channels}
    \item Pacotes Flutter Community (\emph{plus}): \url{https://github.com/fluttercommunity/plus_plugins}
    \item permission\_handler: \url{https://pub.dev/packages/permission_handler}
    \item flutter\_blue\_plus: \url{https://pub.dev/packages/flutter_blue_plus}
    \item ESP32 com Arduino: \url{https://docs.espressif.com/projects/arduino-esp32}
    \item MQTT: \url{https://mqtt.org} \quad mqtt\_client: \url{https://pub.dev/packages/mqtt_client}
  \end{itemize}
\end{frame}

\begin{frame}{Respostas das questões de revisão}
  \begin{respostas}
    \item O sensor transforma temperatura em tensão; o ADC do ESP32 amostra e quantiza em 12 bits; o \emph{firmware} converte em °C e envia por BLE ou MQTT; o app decodifica e exibe. A taxa de amostragem define as leituras por segundo; acima do necessário, só gasta energia.
    \item Ele mede a gravidade somada à aceleração do movimento e é ruidoso: com o aparelho em movimento, a inclinação calculada erra. A fusão com o giroscópio (vetor de rotação) resolve.
    \item Quatro: três coordenadas e o erro do relógio do receptor. Dentro de prédios, o sinal é bloqueado ou refletido e o sistema recorre a Wi-Fi e torres, menos precisos.
    \item No momento em que o recurso é usado, explicando o motivo quando não for óbvio. Se o usuário negar, o app segue funcionando sem o recurso e oferece um atalho para as configurações.
    \item O rádio fica quase sempre desligado e só acorda nos eventos de conexão. O intervalo de conexão é o tempo entre esses encontros, de 7,5\,ms a 4\,s.
    \item Leituras: QoS 0 ou 1, pois uma perdida logo é substituída pela seguinte. Relé: QoS 1 com comando idempotente (ou QoS 2), pois perder o comando é grave e repeti-lo não pode causar dano.
    \item Um bit de início, 8 de dados, nenhum de paridade e 1 de parada. Com \emph{baud rates} diferentes, o receptor amostra nos instantes errados e lê lixo.
    \item MQTT ou \emph{push} em vez de \emph{polling}, envios agrupados, taxa de amostragem menor, \emph{notify} no BLE, o pacote \cod{workmanager} para tarefas periódicas e nada de GPS contínuo.
  \end{respostas}
\end{frame}

\encerramento{Módulo 14 --- Testes: unidade, widget, \emph{golden} e integração.}

\end{document}
